\documentclass[output=paper,colorlinks,citecolor=brown]{langscibook}
\ChapterDOI{10.5281/zenodo.20541858}



\author{Björn Wiemer\orcid{}\affiliation{Johannes Gutenberg-Universität Mainz}}
\title[{The entanglement of grammaticalization and hypoanalysis}]{The entanglement of grammaticalization and hypoanalysis in the rise of futures in Slavic languages}
\abstract{This article pursues two aims. The first one is to show the interrelation of grammaticalization and hypoanalysis in the evolution of future tenses in Slavic languages. For this purpose, hypoanalysis is introduced as a major mechanism of language change, distinct from grammaticalization, and exemplified by the shift of the present tense of perfective stems to mark perfective future by default, which is characteristic of entire North Slavic. This type of hypoanalysis is set off from futurate presents. The second aim consists in specifying the place of this type of hypoanalysis among the means of future marking which have been attested in Slavic languages over the past 1,000+ years. Against this backdrop, the article details the chronological relationship and the dynamics of the areal distribution of the main sources of future tenses in Slavic languages. The author asks about the reasons for North-South splits between Slavic languages in the domain of futures tenses and forwards proposals in order to understand how, in this grammatical domain, grammaticalization and hypoanalysis might have interfered to yield the contemporary distribution of future tenses in Slavic languages.}

%move the following commands to the "local..." files of the master project when integrating this chapter


% \usepackage[markup=underlined]{changes}

% \usepackage{enumitem}
% \usepackage{hyphenat}
% \def\hyph{-\penalty0\hskip0pt\relax}
% \definechangesauthor[color=NavyBlue]{SH}

\IfFileExists{../localcommands.tex}{
   \addbibresource{../localbibliography.bib}
   \usepackage{tabularx,multicol}
\usepackage{url}
\urlstyle{same}

\usepackage[english]{babel}
 
\usepackage{langsci-optional}
\usepackage{langsci-lgr}
\usepackage{langsci-gb4e}

\usepackage{langsci-branding}

%\usepackage{tabularx}

% elena
\usepackage{graphicx} %package to manage images
\usepackage{paralist}

% sarah & lena
\usepackage[justification=centering]{caption}
\usepackage{rotating}

% robin
\usepackage{lscape}
\usepackage{tikz}
\usetikzlibrary{tikzmark}

% dominika
\usepackage{multirow}
\usepackage{hanging}

% olaf & dylan
\usepackage{subcaption}

% svetlana
\usepackage{verbatim}
\usepackage{enumitem}

% wiemer
%\usepackage[markup=underlined]{changes}
%\usepackage{hyphenat}

% wir
\usepackage{longtable}

% for crossref in the introduction jiabin
\usepackage{xr}
   % \newcommand*{\orcid}{}
%

\makeatletter
\let\thetitle\@title
\let\theauthor\@author
\makeatother

\newcommand{\togglepaper}[1][0]{
   \bibliography{../localbibliography}
   \papernote{\scriptsize\normalfont
     \theauthor.
     \titleTemp.
     To appear in:
     E. Di Tor \& Herr Rausgeberin (ed.).
     Booktitle in localcommands.tex.
     Berlin: Language Science Press. [preliminary page numbering]
   }
   \pagenumbering{roman}
   \setcounter{chapter}{#1}
   \addtocounter{chapter}{-1}
}
%
\newbool{bookcompile}
\booltrue{bookcompile}
\newcommand{\bookorchapter}[2]{\ifbool{bookcompile}{#1}{#2}}
%
% %\newcommand{\orcid}[1]{}
\graphicspath{ {./figures/} }
%
\let\eachwordone=\itshape


% Seongha
% \newcommand{\koreex}[1]{{\fontspec{Baekmuk Batang}{#1}}}

% Kutida
\newfontfamily\thaifont{Noto Sans Thai} 
\newcommand{\textthai}[1]{{\thaifont #1}} % für Thai


\newfontfamily\chinesefont{Noto Sans CJK SC}
\newcommand{\textchinese}[1]{{\chinesefont #1}} % Für Chinesisch

% wiemer
\def\hyph{-\penalty0\hskip0pt\relax}
% \definechangesauthor[color=NavyBlue]{SH}



\newcolumntype{K}{>{\raggedright\arraybackslash}p{0.6\textwidth}}  % Right column with 0.7 width


% \newfontfamily\krn[Scale=MatchLowercase,BoldFont=SourceHanSerif-Bold.otf]{SourceHanSerif-Regular.otf}
% \AdditionalFontImprint{Source Han Serif KO}
% \XeTeXlinebreaklocale 'ko'
% \renewcommand{\krn}{}



% Rhee
\newfontfamily\koreanfont{Noto Sans CJK KR} 
\newcommand{\textkorean}[1]{{\koreanfont #1}} % für Koreanisch

\renewcommand{\REF}[2][]{(\ref{#2}#1)}
   %% hyphenation points for line breaks
%% Normally, automatic hyphenation in LaTeX is very good
%% If a word is mis-hyphenated, add it to this file
%%
%% add information to TeX file before \begin{document} with:
%% %% hyphenation points for line breaks
%% Normally, automatic hyphenation in LaTeX is very good
%% If a word is mis-hyphenated, add it to this file
%%
%% add information to TeX file before \begin{document} with:
%% %% hyphenation points for line breaks
%% Normally, automatic hyphenation in LaTeX is very good
%% If a word is mis-hyphenated, add it to this file
%%
%% add information to TeX file before \begin{document} with:
%% \include{localhyphenation}

\hyphenation{
    con-tri-bute
    par-a-digm
    pe-ri-phra-ses
    achieve-ment
    accomp-lish-ment
    We-li-nchen-kang-ci-kok
    cross-ling-u-is-tic
    rec-ord-s
    trun-ca-ted
    de-ri-va-tion-al
    stem-de-ri-va-tion-al
    pre-sent
    in-flu-en-ced
    sys-tem
    ques-tions
    Bor-kovs-kij
    Da-ny-len-ko
    e-qui-va-lents
    re-cog-nized
    In-ter-fe-renz-schicht
    Aus-ei-nan-der-se-tz-ung
    Früh-neu-hoch-deutsch-kor-pus
    Ak-tions-art
    Smir-no-va
    Whi-le
    Ar-chive
    bi-ginnan
    Rei-chen-bach
    ger-ma-nis-ti-sche
    Re-fe-renz-kor-pus
    le-xi-ka-lisch-en
    Alt-hoch-deutsche
    Li-te-ra-tur
    dia-chron
    Dia-lek-to-lo-gie
}

\hyphenation{
    con-tri-bute
    par-a-digm
    pe-ri-phra-ses
    achieve-ment
    accomp-lish-ment
    We-li-nchen-kang-ci-kok
    cross-ling-u-is-tic
    rec-ord-s
    trun-ca-ted
    de-ri-va-tion-al
    stem-de-ri-va-tion-al
    pre-sent
    in-flu-en-ced
    sys-tem
    ques-tions
    Bor-kovs-kij
    Da-ny-len-ko
    e-qui-va-lents
    re-cog-nized
    In-ter-fe-renz-schicht
    Aus-ei-nan-der-se-tz-ung
    Früh-neu-hoch-deutsch-kor-pus
    Ak-tions-art
    Smir-no-va
    Whi-le
    Ar-chive
    bi-ginnan
    Rei-chen-bach
    ger-ma-nis-ti-sche
    Re-fe-renz-kor-pus
    le-xi-ka-lisch-en
    Alt-hoch-deutsche
    Li-te-ra-tur
    dia-chron
    Dia-lek-to-lo-gie
}

\hyphenation{
    con-tri-bute
    par-a-digm
    pe-ri-phra-ses
    achieve-ment
    accomp-lish-ment
    We-li-nchen-kang-ci-kok
    cross-ling-u-is-tic
    rec-ord-s
    trun-ca-ted
    de-ri-va-tion-al
    stem-de-ri-va-tion-al
    pre-sent
    in-flu-en-ced
    sys-tem
    ques-tions
    Bor-kovs-kij
    Da-ny-len-ko
    e-qui-va-lents
    re-cog-nized
    In-ter-fe-renz-schicht
    Aus-ei-nan-der-se-tz-ung
    Früh-neu-hoch-deutsch-kor-pus
    Ak-tions-art
    Smir-no-va
    Whi-le
    Ar-chive
    bi-ginnan
    Rei-chen-bach
    ger-ma-nis-ti-sche
    Re-fe-renz-kor-pus
    le-xi-ka-lisch-en
    Alt-hoch-deutsche
    Li-te-ra-tur
    dia-chron
    Dia-lek-to-lo-gie
}
   \boolfalse{bookcompile}
   \togglepaper[8]%%chapternumber
}{}

\begin{document}
\maketitle

\section{Introduction} \label{sec:wiemer:1}
The main purpose of this article is to show and, to some extent, explain the interrelation of grammaticalization and hypoanalysis in the evolution of future tenses in Slavic languages. In the first place, this presupposes that grammaticalization and hypoanalysis are distinguished as two separate mechanisms of change, that is, they must be disentangled conceptually, before we can demonstrate how they get entangled in observable or reconstructible diachronic processes. I will argue that the comparative study of future tenses in Slavic languages, both in their evolution and their contemporary stage, supplies excellent ground for a demonstration of grammaticalization and hypoanalysis. In connection with this, the present contribution makes an attempt at giving a comprehensive survey of the conceptual and empirical issues, based on the facts of Slavic languages, that have to be considered when trying to disentangle the often entangled relation between grammaticalization and hypoanalysis.

To achieve this aim, I first provide some general information on future tenses in Slavic languages against the backdrop of their inherited tense-aspect system. Concomitantly, I will point out the issues that have to be addressed in order to comply with the main purpose mentioned above and to disclose some diachronic and areal patterns characteristic of Slavic languages with regard to future tenses (Section \ref{sec:wiemer:2}). These issues will be specified at the end of Section \ref{sec:wiemer:2}.

A few words are due concerning the standard divisions among Slavic languages: ‘North Slavic’ is the conjunction of East and West Slavic, that is Slavic spoken north of the Danube and the Pannonian Plain. ‘South Slavic’ is the remainder, roughly south of these topographic points of orientation, with the exception of Serbian in the Vojvodina (spoken north and east of the Danube). South Slavic further divides into Balkan Slavic in the southeast (= Bulgarian, Macedonian, SE-Serbian dialects), Slovene in the northwest, and, in between, Central South Slavic\footnote{This term was suggested by a reviewer, which I hereby gratefully take over.} (CentSS), which corresponds to BCMS, i.e. SerBoCroatian incl. Montenegrin. CentSS subdivides into Štokavian in the east, which also forms the dialectal basis of Standard Serbian and Croatian, Čakavian on the Dalmatian coast, and Kajkavian in the northwest, which comprises entire Slovenian (as the basis of its standard variety) and a small region of adjacent Croatian dialects.

One further “technical” remark: extensive glossing will be supplied only if required for the particular point being made. Otherwise, I will just indicate aspect and tense with superscript indexes (\textsuperscript{\textsc{ipfv}},\textsuperscript{ \textsc{pfv}}).

\section{General background: setting the stage}\label{sec:wiemer:2}

Like predecessors of modern Germanic languages, Common Slavic (3\textsuperscript{rd}--8\textsuperscript{th} centuries AD)\footnote{This term relates to the late Proto-Slavic period (with an assumed still relatively homogeneous dialect continuum) and can be considered equivalent to proto-stages of other IE (sub)branches (cf. \cite{Andersen1985}, \citeyear{andersen_2006b}; \cite{TownsendJanda1996}).} lacked a dedicated future tense (henceforth ‘future’). That is, there was no regular form, or construction, employed as a predicate whose core function was the denotation of a single (‘episodic’) situation located in time after a reference interval (either the speech interval, in deictic tense use, or some other anchor, e.g. in past-tense narratives).\footnote{This more precise understanding of ‘future’ excludes transpositional use of present tense (as with futurate presents, see Section \ref{sec:wiemer:3.1.3}), but also reference to future events by some generalized non-past (as, e.g., in Finnish).} In Common Slavic (CS), apart from a system of perfects (based on a \textsc{be}-copula and an active anteriority participle suffixed with \babelhyphen{nobreak}\textit{l}\babelhyphen{nobreak}), there only existed a contrast between a past domain (with an aorist-imperfect distinction) and an undifferentiated non-past: only present tense stems were used to refer to events posterior to speech time (as for an inherited future perfect see Section \ref{sec:wiemer:3.2.1}). Various periphrastic ways of marking future were in an embryonic stage and could be best considered minor patterns. Possibly already in the CS period, i.e. in times prior to documentation, these patterns were distributed with different biases within the Slavic territory (cf. \cite{andersen_2006b} for an idea of what this distribution might have looked like), but one can hardly speak of any kind of grammaticalized stage (\textit{pace} \cite{andersen_2006b}).

However, already by CS times a perfective–imperfective (pfv.–ipfv.) distinction (a.k.a. perfectivity opposition) had started developing. This opposition is based on stem derivation, consequently pfv. vs ipfv. aspect can be determined in every (finite or non-finite) verb form (see Section \ref{sec:wiemer:3.1}). Over time, this overarching aspect system started interacting with other categories relevant for verb morphology and clausal syntax (\cite{wiemer_serzant_2017}, \cite{wiemer_2020}: 267--268; \citeyear{wiemer_2022}: 70--86). In particular, the strengthening of the perfectivity opposition began to exert its impact on a differentiation of the non-past domain early since it considerably predated the rise of future constructions whose nuclei (auxiliaries) derive from specific lexical sources. These sources include verbs meaning \textsc{become} or \textsc{want}, but also \textsc{begin} (mainly in East Slavic), \textsc{take} (> \textsc{begin}, Ukrainian), \textsc{stand up} (Russ. \textit{stan}-), and most recently the isolated case of \textsc{go} in Slovak (see below); moreover \textsc{have}, prominent at the earliest attested stages, continues playing an important role in Bulgarian and Macedonian (see Section \ref{sec:wiemer:3.2.2}).

Thus, on the one hand, futures that have emerged from grammaticalization are comparatively recent in Slavic languages, and this manifests itself in considerable variation of their various stages over the last 1,000+ years. On the other hand, even a superficial glance at the last centuries (approximately since 1500) reveals a basic North-South split concerning the predominating source domains of morphemes that have turned into nuclei of future constructions: while in the Northern area auxiliaries originating in an inchoative verb (CS *\textit{bǫd-}, \textsc{become}) have come to prevail, the southern area is dominated by clitic auxiliaries deriving from the volition verb \textit{xotěti/xătěti} (\textsc{want}). Slovene is an exception, as it marks the future with a \textsc{become}-auxiliary (\textit{bô}- < *\textit{bǫd-}). Otherwise, however, Slovene “behaves” like any other South Slavic language (and markedly differs from North Slavic), namely: in all South Slavic languages, future auxiliaries can be combined with both pfv. and ipfv. stems; this lack of restrictions as for aspect entails a clear morphological contrast between present and future not only for ipfv., but also for pfv. stems. By contrast, in North Slavic languages future auxiliaries (\textsc{become} or other, minor ones) do not combine with pfv. stems (\cite{wiemer_2020}: 266--270, 275--279). Exceptions are known for colloquial and dialectal varieties of Sorbian (cf. \cite{Schuster-Šewc1996}: 169, \cite{Scholze2008}: 253--254, \cite{Breu2012}: 251--252), and very few instances appear to be attested for historical stages of other North Slavic varieties. Another exception is the Slovak ‘prospective’ construction with a \textsc{go}-verb, which probably appeared in the 18\textsuperscript{th} century and at present shows signs of developing into a future. In this construction pfv. stems can also appear, and this can be explained by the semantics of this construction (cf. \cite{giger_2008}).

These exceptions do not change the general picture: North Slavic languages use the morphological present tense of pfv. stems to mark pfv. future, which acts as a default interpretation of their present tense inflection. There is no additional marker, and the process responsible for this change may be captured under “hypoanalysis”: peripheral and contextually conditioned functions of a construction, or a categorial opposition, turn into its inherent property due to a shift of prominence within an inherited function range. This shift in default meaning distinguishes hypoanalyzed future from futurate presents (see Section \ref{sec:wiemer:3.1.3}). Simultaneously, this shift cannot be subsumed under any of the current understandings of secondary grammaticalization:\footnote{For a critical overview of phenomena that have been called ‘secondary grammaticalization’ cf. \citet{breban_2014}.} the present>future shift is conditioned by the interplay of an age-honoured present tense morphology and a more recent, though likewise well-established stem\hyp{}derivational opposition of pfv. and ipfv. aspect (see Section \ref{sec:wiemer:3.1.1}), but the latter has by no means resulted from any usual grammaticalization process (\cite{wiemer_2020}: 269--270; \citeyear{wiemer_2022}: 67--90).

At least against a European background, the interaction of futures based on the grammaticalization of lexical source expressions with the pfv. future resulting from hypoanalysis looks unique. This justifies differentiating between grammaticalization and hypoanalysis as two distinct types of grammatical change, and the diachrony of future marking in Slavic languages can be considered a paradigm example for the study of their interaction. 

Accordingly, this article pursues two aims. A more general goal consists in outlining the diachronic relation between futures arisen from grammaticalization and from hypoanalysis in Slavic languages; this implies asking for the reasons for the aforementioned North-South split within Slavic. Among the now grammaticalized futures, I will focus on the predominant “players” based on \textsc{want} (South) and \textsc{become} (North), while the others take a backseat. 

The second topic pursued is more specific and can, as it were, be located inside this general picture. The basic North-South split for constructions marking future started gaining salience during the 13\textsuperscript{th}--15\textsuperscript{th} centuries. This follows from the available data and insights gained so far concerning the role of *\textit{bǫd}- ‘become’ in North Slavic and of \textit{xotěti} ‘want’ in South Slavic (with the NW-corner of the latter as an intermediate area); see Section \ref{sec:wiemer:3}. Upon closer inspection, these insights correlate with the behaviour of certain verbal prefixes in the languages at the western edges of North and South Slavic, which can be analyzed as future markers on the present tense of ipfv. stems of certain verb lexemes. Compare \xref{ex:wiemer:1}, where square brackets indicate that the prefix \textit{po}- takes scope over the entire present tense form of the ipfv. stem (infinitive \textit{sunout se} ‘move’) and thereby causes its future interpretation; by contrast, with the pfv. stem \textit{objevit se} ‘appear’ it is the present tense itself which moves the reference to a future event:

\ea Czech \label{ex:wiemer:1}\\
    \gll \textit{Objeví} \textit{se} \textit{mrak} \textit{a} \textit{celé} \textit{tři} \textit{hodiny} \textit{\textbf{se}} \textit{\textbf{po-sun-e}} \textit{obloh-ou}.\\
    appear.\textsc{pfv-(prs>fut.3sg)} \textsc{refl} cloud-\textsc{(nom.sg)} and entire three hours \textsc{refl} \textsc{pfx>fut-}[move.\textsc{ipfv-prs.3sg}] firmament-\textsc{ins}\\
    \glt `A cloud \textbf{will} appear and \textbf{move} across the sky for three hours.’\\
    \hfill (from \cite{ŠevčíkováPanevová2018}: 185)
\z

While this phenomenon has been known for quite some time and described reasonably well, primarily for contemporary Czech, a couple of unsolved questions have remained, which I will identify in Section \ref{sec:wiemer:4}. First of all, these questions should be considered together with differences in the gradual strengthening of the stem\hyp{}derivational aspect opposition and the split between South and North Slavic concerning the interaction of non-past tense with auxiliaries as explicit markers of future. Most briefly, I claim that prefixation as a means to mark future of ipfv. present tense forms was caused by actional lability among certain semantic classes of stems which was more typical of earlier stages in the differentiation among Slavic languages. That is, in the 13\textsuperscript{th}--15\textsuperscript{th} centuries the relevant stems were not yet firmly assigned to pfv. or ipfv. aspect, and the opposition between a class of pfv. and a class of ipfv. stems was less stable on the whole. In addition, future as a function of prefixes added to ipfv. present tense stems resulted from ousting alternative functions that are regularly associated with PFV.PRS. More or less at the same time, we observe the onset of a major split between North and South Slavic concerning the predominant source expressions of future auxiliaries (see above).

Correspondingly, the remainder of this article is structured as follows. In Section \ref{sec:wiemer:3}, I supply the necessary background concerning the main future marking devices, namely the default interpretation of PFV.PRS as (pfv.) future in North Slavic and its preconditions (Section \ref{sec:wiemer:3.1}), and the most prominent constructions based on auxiliaries, namely on \textsc{become} and \textsc{want} (Section \ref{sec:wiemer:3.2}). In Section \ref{sec:wiemer:4}, I concentrate on future marked by prefixes on ipfv. present tense stems (mainly in Czech). In Section \ref{sec:wiemer:5}, both threads are brought together, before I provide a summary and tentative conclusions in Section \ref{sec:wiemer:6}.



\section{Future marking in Slavic}\label{sec:wiemer:3}
The onset of the stem\hyp{}derivational aspect system clearly predates historical records. We can confidently locate its emergent stages in Common Slavic (CS), which overlapped with the time of the Great Migrations (3\textsuperscript{rd}--7\textsuperscript{th} centuries AD). In Old Church Slavonic (since late 9\textsuperscript{th} century AD) we already encounter the basic architecture of the contemporary system (\cite{EckhoffJanda2014}, \cite{eckhoff_2020}, \cite{Kamphuis2020}), although it was less consistent in terms of morphological patterns (combining prefixes and suffixes) and of the functional distribution between related stems. Moreover, morphological patterns and functional distribution had still less consistently spread over different actionality classes (cf. \cite{wiemer_serzant_2017}: 260--272, with further references). The particular relevance of this last point will become evident in Section \ref{sec:wiemer:4}. Overt future markers, to the contrary, started emerging on the basis of an already well-established stem\hyp{}derivational PFV:IPFV opposition. CS did not have any dedicated future, and the earliest attested stages of Slavic languages display a variety of periphrastic constructions used to mark future events, which were based on etyma with the meanings \textsc{become} (> \textsc{be}), \textsc{have}, \textsc{want}, \textsc{begin}, \textsc{take} (> \textsc{begin}); cf. \citet{Birnbaum1958}, \citet{Křížková1960}, \citet[174--185]{Večerka1993}, \citet[91--93]{Grković-Mejdžor2012}, also \citet[146--157]{Borkovskij1949} and \citet[]{seveleva2021} on East Slavic, as for \textsc{take} cf. \citet[18--26]{Danylenko2012} on Ukrainian and \citet[31]{andersen_2006b} on North Russian.

I will begin with the premises that made hypoanalysis possible in North Slavic (Section \ref{sec:wiemer:3.1}). Subsequently, I will dwell upon the two most prominent construction types, \textsc{become} for North Slavic, \textsc{want} for South Slavic, with some more background (Section \ref{sec:wiemer:3.2}).



\subsection{From perfective present to perfective future}\label{sec:wiemer:3.1}
In order to understand the conditions for the pfv. present > pfv. future hypoanalysis in North Slavic we need to spell out a few basics about the architecture of the stem\hyp{}derivational aspect opposition (Section \ref{sec:wiemer:3.1.1}) and its interaction with non-past tense (Section \ref{sec:wiemer:3.1.2}), before we turn to hypoanalysis (Section \ref{sec:wiemer:3.1.3}). 

In general, I take pfv. aspect as a grammatical property by which a verb form indicates that the denoted situation is bounded (limited). Importantly, this need not imply the attainment of some inherent goal (as with telic verbs or clausal configurations); instead, boundedness may simply mark atelic situations as temporally limited (e.g., Russ. \textit{Oni polčasa }[a] \textit{posideli}\textsuperscript{\textsc{pfv}}\textit{ na skamejke, potom }[b] \textit{vstali}\textsuperscript{\textsc{pfv}}\textit{, }[c]\textit{ poguljali}\textsuperscript{\textsc{pfv}}\textit{ i }[d] \textit{rozošlis’}\textsuperscript{\textsc{pfv}} ‘They [a] sat on a bench for half an hour, then [b] got up, [c] walked (for a while) and [d] separated’).


\subsubsection{From prolific stem derivation to grammatical complementarity}\label{sec:wiemer:3.1.1}
In Slavic, aspect is not indicated by unequivocal, monofunctional morphemes. Instead, the morphology of aspect builds on the functional reinterpretation of patterns of stem derivation; these patterns consist in the extension of simpler stems by the (usually consecutive) addition of prefixes and suffixes, or in a suffix change (\cite{breu_2000}, \cite{wiemer_2008}, \cite{wiemer_serzant_2017}, among many others). Apart from minor, mostly lexically restricted or obsolete patterns, two patterns predominate and are productive across contemporary Slavic. See (\ref{ex:wiemer:2}--\ref{ex:wiemer:3}) with infinitives in Polish. The concrete shape of the suffixes is of minor concern (see fn. \ref{fn:5}); stems are given in square brackets.


\ea \footnotesize Predominant patterns of stem derivation and their grammatical reinterpretation \label{ex:wiemer:2}
    \ea \label{ex:wiemer:2a}
    \footnotesize simplex $\Rightarrow$ \textsc{prefix} + simplex\\

    \ex \label{ex:wiemer:2b}
    \hspace{3.5em} \footnotesize prefixed stem $\Rightarrow$ [prefixed stem] + \textsc{suffix} \\
    \hspace{7em} $\Downarrow$ \footnotesize reinterpretation: (a) (potentially) identical lexical concept \\
    \hspace{14.7em} \footnotesize (b) different grammatical distribution\\
    \z
\z

\newpage

\ea \label{ex:wiemer:3}
    \ea \label{ex:wiemer:3a}
    \footnotesize simplex\textsuperscript{\textsc{ipfv}} \rule{0.5cm}{0.1mm} prefixed stem\textsuperscript{\textsc{pfv}} (e.g., Pol. [\textit{gotowa}]-ć $\Rightarrow$ [\textbf{\textit{u}}-\textit{gotowa}]-\textit{ć} ‘cook’)\\
    
    \ex \label{ex:wiemer:3b}
    \hspace{5.2em} \footnotesize prefixed stem\textsuperscript{\textsc{pfv}} \rule{0.5cm}{0.1mm} [[prefixed stem] + \textsc{suffix}]\textsuperscript{\textsc{ipfv}}\\
    \hspace{6.6em}\scriptsize e.g., Pol. [\textit{prze}-\textit{kaz-a}]-\textit{ć} $\Rightarrow$ [\textit{prze-kaz-}\textbf{\textit{ywa}}]-\textit{ć} ‘convey’ (suffix addition\footnote{Most probably, the suffix -\textit{yva}- is an accretion of an earlier suffix -\textit{va}- with a preceding vowel (associated with causatives). On a synchronic level, the relation between (prefixed) pfv. stem and its suffixed ipfv. derivative looks like a suffix change (here: with the thematic vowel -\textit{a}-). For the consistency of the pattern itself this is of minor concern.\label{fn:5}}) \\
    \hspace{7.4em}\scriptsize [\textit{s-praw-i}]-\textit{ć} $\Rightarrow$ [\textit{s-prawi}-\textbf{\textit{a}}]-\textit{ć} ‘cause’ (suffix change) \\
    \z
\z

In a nutshell, the reinterpretation consists in the derivation of more complex stems from simpler stems which can (but need not) share an identical lexical meaning but in any case differ in their distribution over grammatically or pragmatically definable contexts and in their sensitivity to actionality (events, processes, states) and external pluractionality (regular or irregular iteration of distinct eventuality tokens). Grammatical contexts may be defined by the interaction of stems with other categories marked on the verb, the verb phrase or on clause level (such as tense, voice, or polarity); alternatively, there may be pragmatic contexts to which aspect choice is sensitive, such as illocutionary contrasts (e.g., warning vs prohibition) or presuppositions that the speaker assumes to be shared (or not) with the addressee. Importantly, morphemes marking tense, mood, agreement (person/number), but also non-finite forms, all occur outside of stems.\footnote{Present, or non-past, tense can also be indicated by the choice of the stem (compare \ref{ex:wiemer:4} vs \ref{ex:wiemer:5}). However, distinctions between present/non-past stem and infinitive (< aorist) stem are older than the prefixation-suffixation patterns in (\ref{ex:wiemer:2}--\ref{ex:wiemer:3}).} Thus, a form like Pol. \textit{pisała} in \xref{ex:wiemer:4a} has the same structure of tense-agreement marking as has its pfv. counterpart \textit{napisała} in \xref{ex:wiemer:4b} derived with the prefix \textit{na}-.

\ea Polish, past tense \label{ex:wiemer:4}\\
    \ea \label{ex:wiemer:4a}
        \gll [\textit{pis-a}]\textsc{\textsuperscript{ipfv}}\textit{-ł-a}\\
             write-\textsc{thv-pst-sg.f}\\
        \glt `She wrote\textsc{\textsuperscript{ipfv}}’\\
    
    \ex \label{ex:wiemer:4b}
        \gll [\textit{\textbf{na}-pis-a}]\textsc{\textsuperscript{pfv}}\textit{-ł-a} \textit{(list)}\\
             \textsc{\textbf{pfx}}-write-\textsc{thv-pst-sg.f} letter-(\textsc{acc.sg})\\
        \glt `She wrote\textsc{\textsuperscript{pfv}} (a letter)’\\
    \z
\z

The same applies to non-past stems, as in Pol. \textit{roz-wiąż-ą} \xref{ex:wiemer:5a} vs \textit{roz-wiąz-\textbf{uj}-ą} \xref{ex:wiemer:5b}, or the infinitives \textit{prze-kaza-ć }$\Rightarrow$\textit{ prze-kaz-\textbf{ywa}-ć }‘convey’ in \xref{ex:wiemer:3b}, in which a suffix is used to derive an ipfv. stem:

\ea Polish, present tense \label{ex:wiemer:5}\\
    \ea \label{ex:wiemer:5a}
        \gll [\textit{roz-wiąż}]\textsc{\textsuperscript{pfv}} \textit{-ą} \textit{(sznurk-i)}\\
             \textsc{pfx}-tie- \textsc{prs.3pl} lace-\textsc{acc.pl}\\
        \glt `They (will) untie (the shoelaces).’\\
    
    \ex \label{ex:wiemer:5b}
        \gll {[\textit{roz-wiąz-\textbf{uj}}]{\textsc{\textsuperscript{ipfv}}}} {\textit{-ą}} \textit{(sznurk-i)}\\
             {\textsc{pfx}-tie-\textsc{\textbf{sfx}}} {\textsc{-prs.3pl}} {lace-\textsc{acc.pl}}\\
        \glt `They untie / are untying (the shoelaces).’\\
    \z
\z

Present tense stems, as in \xref{ex:wiemer:5a} and \xref{ex:wiemer:5b}, often show morphonological alternations at the end of roots; compare -\textit{wiąż}- [v'õʒ] in \xref{ex:wiemer:5a} vs -\textit{wiąz}- [v'õz] in \xref{ex:wiemer:5b}. Likewise, suffixes may alternate between the infinitive and the present (or non-past) stem; compare -\textit{ywa}- (see \textit{przekaz-\textbf{ywa}-ć} in \ref{ex:wiemer:3b}) vs -\textit{uj}- (see \textit{rozwiąz-\textbf{uj}-ą} in \ref{ex:wiemer:5b}). These processes affect inflectional morphemes only insofar as present tense stems imply a specific set of person/number endings.\footnote{This principle is an IE heritage. One may thus say that present tense is indicated twice: by the stem and by its correlated endings.} Otherwise stem derivation does not affect the explicit marking of the respective categories, nor is it itself affected by them. However, it often restricts combinability with these categories or their interpretation. This is what happens, in particular, to tense distinctions and the combinability of the infinitive with all kinds of auxiliary verbs. Thus, forms like \textit{rozwiążą} ‘they (will) untie’ in \xref{ex:wiemer:4a}, which are based on the non-past stem, are by default interpreted as future.\footnote{Strictly speaking, such forms might better be glossed as non-past (\textsc{npst}) – first of all, because non-future interpretations have not been ousted (see Section \ref{sec:wiemer:3.1.2}).}

Suffixation techniques had been part and parcel of stem derivation since Proto-Indo-European (PIE) times, while more recent prefixation developed from morphologically independent verbal modifiers and became commonplace in practically all IE groups of Europe. However, only Slavic brought prefixation beyond the stage of merely “telicizing bounders” and applied it all over the verbal lexicon in combination with suffixes.\footnote{\label{fn:9}Practically all other IE languages of Europe lost suffixes as a technique to mark actionality distinctions independently of tense and valency. Lithuanian is an exception, possibly via contact with (East) Slavic (\cite{Arkadiev2014}; \citeyear{Arkadiev2015}, \cite{wiemer_serzant_2017}: Section \ref{sec:wiemer:5}, \cite{wiemer_2023}: 205--212).} Since both strategies started to combine (no later than in CS), they “stayed together” as prolific devices that continued to be productive through the entire history of Slavic languages, albeit to varying degrees in different subareas.

As a consequence of prolific stem derivation, patterns like those in (\ref{ex:wiemer:2}--\ref{ex:wiemer:3}) have been creating a tremendous number of morphologically related verb stems. In many cases, the derivational affixes did not change the lexical meaning, but yielded synonyms, which started to be distributed over different grammatical contexts, with an increasing tendency towards complementary distribution. In the Slavist tradition, pairs of morphologically related pfv. and ipfv. stems that are able to convey an identical lexical meaning but occur in complementary contexts\footnote{The conditions and tests on aspectual pairedness are concisely reviewed in \citet[227--228]{wiemer_2017}, \citet[236--240]{WiemerWyroślak2020}, and \citet[27--29]{Łaziński2020}.} are called aspect pairs. These can, in a sense, be considered the backbone of the system. 

However, the aspect membership of a particular stem does not hinge on entering into a relationship of strict replacement conditions or on looser relations, e.g. between stems denoting atelic activities (e.g., Russ. ipfv. \textit{plakat’} ‘weep’, \textit{letet’} ‘fly’) and their ingressive modifications (pfv. \textit{za-plakat’} ‘start weeping’ and \textit{po-letet’} ‘fly off’, respectively), which can be considered ‘aspect partnerships’ (\cite{Lehmann1988}). Ultimately, membership in pfv. or ipfv. aspect depends on whether the respective stem adds a meaning of boundedness (limitation) to the context, or is compatible with this feature ($\Rightarrow$ pfv.), or whether it does not, or is indifferent in this respect ($\Rightarrow$ ipfv.). This can be tested in various ways, albeit to different degrees of reliability. For instance, a reliable indicator of ipfv. aspect is the ability of the finite verb to denote ongoing processes (as in \textit{Right now, X is eating his porridge}); however, not all ipfv. stems are capable of denoting ongoing processes.

After all, the spread of derivational patterns like those in (\ref{ex:wiemer:2}--\ref{ex:wiemer:3}) has turned into a productive technique that opposes stems which are morphologically related but show complementary distribution over actional and pluractional contexts (and thence to more and more contrasts). This was a long-lasting process, which only stepwise involved different actionality groups with varying speed and consistency. There are some indications that aspect pairs and partnerships appeared first for stems that coded strictly goal-directed (‘strongly telic’) situations and for some verbs denoting speech acts; only later did pairs for punctual events and for atelic activities appear.\footnote{Cf. \citet{Bermel1997}, \citet[]{mende_1999}, \citet[]{EckhoffJanda2014}, \citet[]{eckhoff_2020}, \citet{Kamphuis2020}. On the significance of perfectivized atelic stems for the aspect system cf. \citeauthor{Arkadiev2014} (\citeyear{Arkadiev2014}; \citeyear{Arkadiev2015}: 82--89) and \citet[265--269]{wiemer_serzant_2017}.} For some of the latter aspect partnerships proved unstable; in particular, this applies to verbs of (directed) motion, as we will see in Sections \ref{sec:wiemer:4.3}--\ref{sec:wiemer:4.4}.


\subsubsection{Irrealis meanings and the conventionalization of invited implicatures}\label{sec:wiemer:3.1.2}

Old Church Slavonic (OCS) presents us with the earliest Slavic text monuments (New Testament, gospels) that have been passed down to us from the 10\textsuperscript{th} century onwards. They already show a fairly well-developed distinction between morphologically related stems with a tendency toward complementary division of labour of actionality distinctions, which can be discerned primarily in narrative past tense contexts. As for the non-past domain, matters might seem less clear inasmuch as there was no dedicated future marking. However, provided pfv. and ipfv. stems can be sufficiently distinguished on the basis of their stem\hyp{}derivational relation, non-past forms of pfv. stems are also encountered in pluractional (i.e. habitual or gnomic) contexts. In such contexts they do not only signal ‘that the event occurs on more occasions, but rather that the event will occur when the right circumstances occur’ (\cite[162]{Kamphuis2020}; emphasis added); see similarly in \citet[152]{Večerka1993}. This amounts to dispositional or circumstantial, i.e. non-deontic modal, readings. Compare, for instance, \xref{ex:wiemer:6}, in which PFV.PRS translates Greek present tense forms, and \xref{ex:wiemer:7}, which contains a relative clause with a generic head noun (\textit{vьsěkъ }‘anybody’).\footnote{\textit{vъzrěti} ‘look’ (ipfv. \textit{zrěti}) translates a Greek present active participle.} In either case, the non-past forms do not refer to any single, let alone ongoing situation, and the named events do not (or not necessarily) occur in some interval posterior to the time of utterance; instead, time reference is generalized, or unspecific:

\ea \label{ex:wiemer:6}
    \textit{i} \textit{g[lago]ljǫ} . \textit{semu} \textit{idi}\textsc{\textsuperscript{ipfv.imp}} . \textit{i} \textit{idetъ}\textsc{\textsuperscript{ipfv.prs}} \textit{i} \textit{drugumu} \textit{pridi}\textsc{\textsuperscript{pfv.imp}} . \textit{i} \textit{\textbf{pridetъ}}\textsc{\textsuperscript{pfv.prs}}. \textit{i} \textit{rabu} \textit{moemu} \textit{sъtvori}\textsc{\textsuperscript{pfv.imp}} \textit{se} \textit{i} \textit{\textbf{sъtvoritъ}}\textsc{\textsuperscript{pfv.prs}}\\
    ‘And I say to this one, “Go!” and he goes, and to another, “Come!” and he \textbf{comes}, and to my slave, “Do this!” and he \textbf{does} it.’\\
    \hfill (Matthew 8:9; \cite{Kamphuis2020}: 162)
\z

\ea \label{ex:wiemer:7}
    \textit{azъ} \textit{že} \textit{g[lago]ljǫ} \textit{vamъ} \textit{ěko} \textit{vьsěkъ} \textit{iže} \textit{\textbf{vъzьritъ}}\textsc{\textsuperscript{pfv.prs}} \textit{na} \textit{ženǫ} \textit{sъ} \textit{poxotijǫ}. \textit{uže} \textit{ljuby} \textit{sъtvori}\textsc{\textsuperscript{pfv.aor}} \textit{sъ} \textit{nejǫ} \textit{vъ} \textit{srъdьci} \textit{svoemь} [Z, M]\\
    ‘But I say to you that everyone who \textbf{looks} at a woman with lust for her has already committed adultery with her in his heart.’\\
    \hfill (Matthew 5:28; \cite{Kamphuis2020}: 162)
\z

However, also well-attested are pfv. non-past forms with an unambiguous future reference, i.e. with reference to single events after (presumed) utterance time; compare, for instance, (\ref{ex:wiemer:8}--\ref{ex:wiemer:9}):

\ea \label{ex:wiemer:8}
    \textit{tъgda} \textit{°gla} \textit{imъ} \textit{°is.} \textit{vsi} \textit{vy} \textit{\textbf{sъblaznite}}\textsc{\textsuperscript{pfv.prs}} \textit{\textbf{sę}} \textit{o} \textit{mně} \textit{vъ} \textit{sьjǫ} \textit{noštь}\\
    ‘Then Jesus told them: all of you (will) \textbf{become concerned} about me this night.’\\
    \hfill (Matthew 26:31; \cite{Večerka1993}: 175)
\z

\ea \label{ex:wiemer:9}
    \textit{azъ} \textit{\textbf{pomlьčǫ}}\textsc{\textsuperscript{pfv.prs}} \textit{sebe} \textit{radi.} \textit{da} \textit{\textbf{povědętъ}}\textsc{\textsuperscript{pfv.prs}} \textit{čudesa} \textit{°ba} \textit{mojego}\\
    ‘I \textbf{will keep silent} myself for them to tell the miracles of my God [or: so that they \textbf{may/will tell} the miracles of my God].’\\
    \hfill (Suprasl’ chronicle; \cite{Ivanova-Mirčeva1962}: 35)
\z

In fact, pfv. non-past was a frequent, if not the most frequent, means to express future events in OCS, also as an equivalent of the Greek future (cf. \cite{Birnbaum1958}: 17--18, \cite{Bunina1959}, \cite{Ivanova-Mirčeva1962}: 28--49, \cite{Večerka1993}: 175--176). Moreover, it was widespread after the irrealis connective \textit{da}, particularly in contexts of volition (e.g., purpose clauses as in \ref{ex:wiemer:9} and \ref{ex:wiemer:10}), and in the protasis of conditionals. The latter belong to cognition-based contexts with downgraded epistemic support, and they are associated with relative clauses with generic referents (as in \xref{ex:wiemer:7}).

\ea \label{ex:wiemer:10}
    \textit{sъlězъ} \textit{otъ} \textit{°nbse} \textit{°ađlъ} \textit{sedmь.} \textit{\textbf{da}} \textit{\textbf{otъženǫtъ}}\textsc{\textsuperscript{pfv.prs}} \textit{vse} \textit{zъlo}\\
    ‘Seven angels came down from heaven \textbf{to drive away} [lit. \textbf{that (they) drive away}] all evil.’\\
    \hfill (from \cite{Večerka1993}: 92)
\z

That is, pfv. non-past was encountered in all kinds of irrealis contexts, i.e. in ‘the realm of the imagined or hypothetical’ to denote a potential event, but ‘not an observable fact of reality’ (\cite{Elliott2000}: 66--67). Future events are closely associated to the realm of irrealis inasmuch as they cannot be checked at the moment of speaking.

For all contexts just mentioned, the [+bounded]-feature of pfv. aspect is crucial, in particular for pluractional contexts (and their non-deontic modal implications): in order to conceive of repetitions, situations need to be discernible at distinct (= limited) intervals, and each potential occurrence of the denoted situation type can easily be linked to conditions (= circumstances) that make it more likely to happen and allow it to represent the situation type as its imagined token. Likewise, if one knows about somebody else’s inclination or about some lawful behaviour (i.e. a situation type), one may -- via abduction -- infer from circumstances observed, or holding, at the moment of speech that a specific (singular) manifestation of this behaviour (i.e. a situation token) will happen. There is thus a natural connection between habituality, the assignment of stable, and thus predictable, properties and predictions based on ‘the right circumstances’ (see quote above) that are observed in a particular setting. These connections have been recognized as a chain of conceptual bridges to future reference (cf. \cite{tatevosov_2004}, \cite{Shluinsky2006}, \cite{Sonnenhauser2008}). The last link of this chain is predictions (i.e. epistemic judgments) concerning specific events. These events cannot be “checked” since they are thought to occur after the speech act (or another reference interval), nor can habitual or associated modal meanings be checked since they are not about specific intervals.

Clearly, the present tense, or rather: non-past, forms of pfv. stems in Slavic were unable to denote any ongoing situations right from the start whenever the perfectivity opposition established itself (during CS) on the basis of the past:non-past contrast in verb morphology. However, given the tight connection between the aforementioned usage types of the non-past forms of pfv. stems in Slavic – and since no further morphology is involved – one wonders how present and future can be told apart for these forms. After all, all sorts of readings as “inactual present” were no less inherent to these forms from the start than reference to the future, i.e. to single events after speech time. Therefore, if future (i.e. reference to a single situation token after speech time) has become predominant for pfv. non-past in North Slavic (see Section \ref{sec:wiemer:3.1.4}), this can only mean that this function has become more salient within this domain. It does not mean that non-future functions (related to irrealis) have vanished; instead, they have only been driven into the background as being less frequent.

\subsubsection{Hypoanalysis}\label{sec:wiemer:3.1.3}

Before we continue browsing through diachrony, let us examine to what extent such a shift of saliency within the non-past can, and should, be subsumed under hypoanalysis. \citet[126--127]{Croft2000} defines hypoanalysis as follows:

\begin{quote}
In \textsc{hypoanalysis}, the listener reanalyzes a contextual semantic/functional property as an inherent property of the syntactic unit. In the reanalysis, the inherent property of the context (often the grammatical context (...)) is then attributed to the syntactic unit, and so the syntactic unit in question gains a new meaning or function.
\end{quote}

In Croft’s wording, hypoanalysis applies to syntactic units and it results in a shift to a new meaning that previously was only implied by favourable types of context. Strictly speaking, the case of North Slavic PFV.PRS acquiring the meaning of (pfv.) future as its default function slightly differs in both respects. First, the new prominent function of PFV.PRS was not really new, but already implied as just one among others that an undifferentiated non-past form could acquire, before and after the perfectivity opposition came into play (see Section \ref{sec:wiemer:3.1.2}). Thus, the shift rather consisted in highlighting a subdomain within the non-past, no matter what caused this conventionalized shift of prominence. Second, and more importantly, we are not dealing with a change applying to syntactic units, nor does the change hinge on any source expressions (etyma) in their syntagmatic settings; instead it pertains to the interpretation of functional oppositions marked, or induced, by patterns of stem derivation. It is, thus, paradigmatic contrasts between two classes of stems (called ‘pfv.’ and ‘ipfv.’) against the background of morphological patterns which favour the default shift into future. Therefore, in subsuming the conventionalized future shift of PFV.PRS under hypoanalysis we should be aware that this brings about a certain expansion of this notion beyond Croft’s original definition.

Moreover, paradigmatic contrasts lie at the core of \citegen{andersen_2006b} considerations about types of language change which he applies to the evolution of futures in Slavic languages. Andersen understands grammaticalization as a complex of changes, more or less in line with mainstream literature on grammaticalization, but he neither explicitly mentions hypoanalysis, nor does he consider the default shift of PFV.PRS from undifferentiated non-past to future.\footnote{Simultaneously, \citet[31]{andersen_2006b} erroneously describes pfv. aspect as referring to states. Of course, this misunderstanding precludes an adequate analysis of what happens to PFV.PRS.} However, hypoanalysis comes close to \citegen[39]{andersen_2006b} ‘regrammation’, which amounts to a ``change by which a grammatical expression through reanalysis is ascribed different grammatical content".

Apart from these caveats, another point is crucial: the conventionalized default shift of PFV.PRS to pfv. future in North Slavic cannot be subsumed under futurate presents. The latter have been well described for Germanic languages, but they are certainly widespread beyond. Futurate presents do not change their morphology, which is also the case for the North Slavic “hypoanalyzed” pfv. future. However, futurate presents acquire future reference not by virtue of a conventionalized default; instead, they are just one, non-dominant reading among many. Futurate use is often difficult to predict, although plans and schedules may trigger it (\cite{Copley2014}); see \xref{ex:wiemer:11}. Moreover, the present tense form can be ambiguous between an ongoing present and a future reading (\cite{Hilpert2008}: 157, 159). This is excluded for North Slavic PFV.PRS; instead, in all (not only North) Slavic languages it is the present tense of ipfv. stems which is widely used for scheduled future, just like the English present progressive or the German present. 

This comparison also shows that futurate presents need not build on a contrast of grammatical aspect nor of present vs future. In English, it is the “non-perfective” member of the aspect opposition which is used for scheduled events (although the Simple Present can be used as well; cf. \cite{Copley2008}; \citeyear{Copley2014}: 72) and there is a clear opposition between present and future morphology. None of these properties holds, e.g., for German (there is no grammatical aspect and the \textit{werden}-future is much less consistently employed than the English \textit{will}-future), in which futurate presents are quite commonplace.

We thus have to distinguish between common shifts of temporal reference of present tense forms that, after all, do not change the default function of these forms, on the one hand, and a conventionalized shift (or narrowing) leading to a new default function, on the other. Compare the following translational equivalents: the Russian ipfv. present (\textit{vstrečaemsja} ‘we meet’) for scheduled events is equivalent to the present tense in German, English, and Bulgarian (see \ref{ex:wiemer:11}); by contrast, PFV.PRS (\textit{vstretimsja} ‘we (will) meet’) is functionally equivalent to the future form in languages like English and Bulgarian, where the present-future contrast is more firmly established than in German (see \ref{ex:wiemer:12}). This corresponds to the default of future reference for PFV.PRS, which is not the case with the futurate interpretation of present tense in \xref{ex:wiemer:11}.

\ea \label{ex:wiemer:11} 
{\scriptsize\textbf{scheduled (planned) events}\\
\begin{tabbing}
    \hspace{1.5cm} \= \hspace{5cm} \= \hspace{4cm} \= \kill
    Russ. \> \textit{My vstrečaemsja zavtra.} \> ipfv. present \\
    Bulg. \> \textit{Sreštame se utre.} \> ipfv. present \\
    Engl. \> \textit{We are meeting / meet tomorrow.} \> present (progressive / simple) \\
    Germ. \> \textit{Wir treffen uns morgen.} \> present (no aspect distinctions) \\
\end{tabbing}
}
\z


\ea \label{ex:wiemer:12}
{\scriptsize \textbf{future (predictive, but not necessarily planned)}\\
\begin{tabbing}
    \hspace{1.5cm} \= \hspace{5cm} \= \hspace{6cm} \= \kill
    Russ. \> \textit{My vstretimsja skoro / kogda-nibud}’. \> pfv. present (non-past) > pfv. future \\
    Bulg. \> \textit{Šte}\textsc{\textsuperscript{fut}} \textit{se sreštnem skoro / njakoj den}. \> pfv. future \\
    Engl. \> \textit{We will}\textsc{\textsuperscript{fut}} \textit{meet soon / some other day.} \> future \\
    Germ. \> \textit{Wir treffen uns bald / ein andermal.} \> present \\
    \> \textit{Wir werden}\textsc{\textsuperscript{fut}} \textit{uns bald treffen.} \> future (optional) \\
\end{tabbing}
}
\z

After these clarifications let us now return to the reconstruction of the “chronology of changes” that contributed to the differentiation of future marking among Slavic languages.


\subsubsection{Beginning split between North and South Slavic}\label{sec:wiemer:3.1.4}

The first written documents of East Slavic date from the 11\textsuperscript{th} century, the first West Slavic documents are from the late 13\textsuperscript{th} century in Old Czech, and the first Old Polish texts are from the mid-14\textsuperscript{th} century. The earliest documents in western varieties of South Slavic are from the 12\textsuperscript{th} century in the Štokavian-Čakavian area (Old Serbian-Croatian),\footnote{\label{fn:14}The oldest texts with clearly Slovene features are the \textit{Freisinger Denkmäler} (about 1000 AD), but they are very limited and do not contribute anything to the issue of (emergent) futures (see Section \ref{sec:wiemer:3.2.1.1}).} whereas it is difficult to name a clear “cut-off point” between Old Church Slavonic and the emergent stages of Balkan Slavic (leading to modern Bulgarian and Macedonian), but Middle Bulgarian is considered to comprise the 12\textsuperscript{th}--14\textsuperscript{th} centuries. Regardless of these different “onsets” of historical attestation, already at these earliest accessible stages some differences can be observed among these groupings (East, West and western vs eastern South Slavic) when it comes to marking future events. As for the emergence of auxiliary-constructions, the employment of \textsc{begin}-verbs (+ infinitive), apart from \textsc{have}, was strongest in early East Slavic (up to the 16\textsuperscript{th}--17\textsuperscript{th} centuries), but unremarkable in the remainder. In turn, in the earliest West Slavic texts we already encounter a clear preference for \textsc{become} (CS *\textit{bǫd}-), which proved a serious “rival” of \textsc{want} and \textsc{have} in the early Štokavian-Čakavian area, before it was superseded by \textsc{want} (see Section \ref{sec:wiemer:3.2}). 

Moreover, in South Slavic dedicated future markers have appeared that are not sensitive to aspect, as they also combine with pfv. stems (see Section \ref{sec:wiemer:3.2.2}). This stands in stark contrast to North Slavic and justifies the assumption that the conventionalization of future meaning as a default of PFV.PRS (= hypoanalysis) was blocked by new constructions with \textsc{begin}- and \textsc{become}-auxiliaries. 

What happened to PFV.PRS? The shortest answer is that in North Slavic the future implicature turned into a conventionalized default, but non-future functions of PFV.PRS have not been abandoned and are still well-attested even in Russian (see \ref{ex:wiemer:13}, with a habitual reading), although with heavier restrictions than in the western parts of North and of South Slavic.

\ea \label{ex:wiemer:13} Russian\\
\textit{I} \textit{s} \textit{ego} \textit{točki} \textit{zrenija,} \textit{novyj} \textit{den’} \textit{--} \textit{ėto} \textit{novoe} \textit{čudo.} \textit{On} \textit{ne} \textit{speša} \textit{\textbf{prosnetsja}}\textsuperscript{\textsc{pfv.prs}}, \textit{plavno} \textit{\textbf{primet}}\textsuperscript{\textsc{pfv.prs}} \textit{vannu,} \textit{vdumčivo} \textit{\textbf{vyp’et}}\textsuperscript{\textsc{pfv.prs}} \textit{kofe,} \textit{(...).} (RNC)\\
‘And from his point of view, a new day is a new miracle. He slowly \textbf{wakes up}, \textbf{takes} a gentle bath, thoughtfully \textbf{drinks} coffee (...).’ \\
\z

By comparison, in contemporary Bulgarian (and Macedonian) PFV.PRS has largely been restricted to co-occurrence with irrealis-marking connectives, mainly with the verbal proclitic \textit{da}, but also with conditional or temporal conjunctions (e.g. Bulg. \textit{ako} ‘if’, \textit{kogato} ‘when’), to relative clauses (especially headless ones, see \ref{ex:wiemer:14}) and to biclausal proverbs with a sequential structure \xref{ex:wiemer:15}. In all such cases, conditionality (‘if..., then...’) keeps these clause units together. In addition, they are deprived of specific temporal reference (see \ref{ex:wiemer:14}--\ref{ex:wiemer:15}).\footnote{For comprehensive accounts of PFV.PRS-uses in Balkan Slavic cf. \citet[231--233]{Maslov1956}, \citet{Lindstedt1985}, \citeauthor{Tomić2006} (\citeyear{Tomić2006}; \citeyear{Tomić2012}).}

\ea \label{ex:wiemer:14} Bulgarian\\
\textit{Kakvoto} \textit{\textbf{reši}}\textsc{\textsuperscript{pfv.prs}}, \textit{go} \textit{pravi}\textsuperscript{\textsc{ipfv.prs}}. \\
‘What(ever) s/he \textbf{decides}, s/he realizes (it).’ \\
\hfill (adapted from \cite{Lindstedt1985}: 192)
\z

\ea \label{ex:wiemer:15} Bulgarian\\
\textit{\textbf{Izlăžeš}}\textsuperscript{\textsc{pfv.prs}} \textit{li} \textit{vednăž,} \textit{ne} \textit{ti} \textit{vjarvat}\textsuperscript{\textsc{ipfv.prs}} \textit{vtori} \textit{păt}. \\
‘If you \textbf{lie} once, people don’t believe you a second time.’ \\
\hfill (from \cite{Maslov1956}: 231)
\z

In all these contexts, PFV.PRS continues to be freely used in Russian and all other North Slavic languages. Therefore, non-future meanings of PFV.PRS are not only original throughout Slavic, but they have stayed the main shared usage domain of all South and North Slavic languages until today. Even in the eastern part of Slavic (both North and South) non-future uses form a strong “residual” of PFV.PRS irrespective of whether dedicated future markers have arisen and how they interact with aspect. That is, neither the default shift of PFV.PRS to pfv. future (hypoanalysis) in North Slavic nor the grammaticalization of future auxiliaries that are “insensitive” to aspect (in South Slavic) have caused any major changes in the irrealis usage domain of PFV.PRS (cf. \citeauthor{Wiemer2026} \citeyear{Wiemer2026} for a detailed discussion). 

Against this backdrop, can the period when this default shift became salient somehow be specified? According to specialists, already by the 12\textsuperscript{th} century North and South Slavic seem to have begun to divide up the preferred domains of usage of PFV.PRS. As for South Slavic, \citet[87]{Grković-Mejdžor2012} remarks that already in the oldest Štokavian/Čakavian documents (12\textsuperscript{th}--15\textsuperscript{th }c.) PFV.PRS exclusively appears in all irrealis domains except the future (cf. also \cite{Belić1999}: 443); evidently, future reference had been taken over by other means (see Section \ref{sec:wiemer:3.2.2}). At first sight, this conclusion seems to contrast with analyses of the evolution of future marking in Balkan Slavic which suggest several centuries of overlap between PFV.PRS and the developing \textsc{want}-future. Thus, following \citet[42--49, 67--68]{Ivanova-Mirčeva1962} or \citet{asenova_guentcheva_2022} one would conclude that, in Bulgarian, PFV.PRS continued to mark future events (also in main clauses) up to the 16\textsuperscript{th} century and, in this usage, ultimately disappeared only by the 18\textsuperscript{th} century. However, such a conclusion would mainly be based on texts which continue the OCS “canon”, whereas texts reflecting vernacular usage largely remained out of focus. Vernacular speech may be assumed to have propelled the newer futures based on \textsc{want} and \textsc{have}. Thus, the impression that these newer constructions and PFV.PRS co-existed over such an amount of time as means of future marking most probably simply reflects an overlay of conservative OCS-traditions with much more innovative vernacular speech (see also Section \ref{sec:wiemer:3.2.2}).\footnote{I am grateful to Jasmina Grković-Mejdžor for having drawn my attention to this fallacy.}

As concerns North Slavic, \citet[145]{Borkovskij1949} points out that during the 12\textsuperscript{th}--16\textsuperscript{th} century PFV.PRS was the predominant choice (88\%) for reference to future events in East Slavic secular texts. Similarly, \citet[124]{misina_2018} summarizes findings on Old East Slavic and notes that in the earliest attested period (11\textsuperscript{th}--14\textsuperscript{th} centuries) the most common way of referring to the future was by using PFV.PRS, that is, just as in OCS (see Section \ref{sec:wiemer:3.1.2}). Provided these conclusions can be generalized for South and North Slavic, respectively, they are indicative of a beginning split between North and South Slavic in the 13\textsuperscript{th}--14\textsuperscript{th} centuries.



\subsection{Grammaticalization: \textsc{become} vs. \textsc{want}}\label{sec:wiemer:3.2}

Caveats concerning discrepancies between vernacular speech and OCS\hyp{}influenced writings are probably appropriate for all constructions that were grammaticalized into futures in Slavic. As a rule, the constructions in which lexical source expressions of emergent future markers and the form of the lexical verb combined were ancient minor patterns. This applies in particular to the combination of \textsc{become} (*\textit{bǫd}-) with the \textit{l}-form (originally an active anterior participle), whereas the combination of \textsc{become} with the infinitive was one of the really few innovations (see Section \ref{sec:wiemer:3.2.1}). In this construction, \textit{l}-form and infinitive must have co-existed for some time in different areas, but data is too scarce and inconclusive, also as for aspect restrictions of the lexical verb and the issue whether and when *\textit{bǫd}- +\textit{ l}-participle turned from a future perfect into a “simple” future. Finally, we also cannot decide to which extent this construction was used in parallel with future reference marked by PFV.PRS (cf. \cite{Křížková1960}: 95--96, 156--158, \cite{Whaley2000}: 26--33).

In South Slavic, this parallelism was probably reduced with the consolidation of the `want'-based future (see Section \ref{sec:wiemer:3.2.2}). However, a complex syntactic and functional distribution of the latter with the \textit{bud}-construction and PFV.PRS continues to be characteristic of different Central South Slavic (CentSS) varieties until today, in which \textit{bud}-constructions are employed like a future perfect in certain embedded clauses.\footnote{Cf. \citet{Kovačević2008} on Serbian, \citet{Vuković2014} for Croatian complex sentences. Cf. also \citet[20]{andersen_2006b} for diatopic variation in the Kajkavian, \citet{Lisac2009} for diatopic variation in the Čakavian region. Cf. also \citet[348--349]{Lukežić2015} for a brief summary on Kajkavian, Čakavian and Štokavian and \citet[11--18]{Grković-Mejdžor2019} for a concise overview of the diachronic development and diatopic distribution of future marking in South Slavic.\label{fn:17}} Regardless of meaning, contemporary Slovene has \textit{bô}- only with the \textit{l}-form. Likewise, adjacent Croatian Kajkavian dialects have \textit{bud}- (or \textit{bô/bu}-) with the \textit{l}-form. For some of them this auxiliary is also occasionally noted with the infinitive (see fn. \ref{fn:17}), but this combination might probably be explained through interference with Štokavian dialects or (likewise Štokavian-based) Standard Croatian.\footnote{\label{fn:18}I am obliged for this suggestion to an anonymous reviewer.} As for North Slavic, only Polish (with contiguous regions) and Kashubian still optionally use the \textit{l}-form (of ipfv. stems) with the future auxiliary (\textit{będ}-),\footnote{\label{fn:19}There are no conceivable functional motivations of this variation in Polish (\cite{Łaziński2006}: 317--321, \cite{blaszczak_etal_2014}: 184--187). In Kashubian the distribution of both variants rather seems to be diatopic, possibly with stronger Polish influence (for \textit{bǫd}- + INF) in the southern part (\cite{andersen_2006b}: 141--142)} in all other North Slavic languages it only occurs with the infinitive. 

\subsubsection{\textsc{become}}\label{sec:wiemer:3.2.1}

The Slavic root *\textit{bǫd}- is cognate to the existential-copular verb \textit{by-ti }‘be’ and can be considered a nasal extension of a PIE root (< PIE *\textit{bhū-n-d(h)}-); cf. \citet{werner_1996}, \citet[95]{Grković-Mejdžor2012}. Probably, *\textit{bǫd}- originally was inchoative,\footnote{Cf. \citet[23--24]{Whaley2000}. For a recent concise overview of positions cf. \citet[570--571]{Pen’kova2019a}.} although by the time of its first written attestations it might have lost its inchoative flavour (\cite{Kraemer2005}: 75--81). Nonetheless, an inchoative feature might have played a role in its reinterpretation as a future marker (see below).

The development of the \textsc{become}-future must have started after the split of Common Slavic into North and South Slavic (i.e. more or less by the 9\textsuperscript{th} century AD), as it took different paths within individual languages (\cite{Whaley2000}: 25, 73). This shows up in two respects: first, whether aspect restricts its use; second, whether *\textit{bǫd}- combines with the \textit{l}-form (originally an active anteriority participle) or the infinitive. Initially, \textit{bǫd}- only combined with the \textit{l}-participle as a future perfect (a.k.a.\textit{ Futur II} or \textit{futurum exactum}), without restrictions of aspect (see \ref{ex:wiemer:16}).

\newpage

\ea Old Church Slavonic \label{ex:wiemer:16}\\
    \gll \textit{i} \textit{ašte} \textit{grěx-y} \textit{\textbf{bǫd-etъ}} \textit{\textbf{sъtvori-l-ъ}} \textit{otdad-ętъ} \textit{sę} \textit{emu}\\
         \textsc{ptc} if sin-\textsc{acc.pl} \textsc{fut-3sg} make[\textsc{pfv}]-\textsc{lpt-sg.m} surrender[\textsc{pfv}]-\textsc{prs.3pl} \textsc{refl} 3\textsc{sg.m.dat}\\
    \glt ‘Even if he commits [lit. \textbf{will have committed}] sins, they will devote themselves to him.’\\
    \hfill (Jak 5:15; \cite{Večerka1993}: 180)
\z

However, already in OCS this temporal function sometimes oscillated with an epistemic, or even inferential, interpretation (\cite{Lépissier1960}, \cite{Ivanova-Mirčeva1962}: 145--147). The subsequent earliest attested stages of Slavic languages already reveal considerable variation.\footnote{For a brief overview cf. \citet[197--199]{arkadiev_wiemer_2020}. Reliable corpus-based surveys and insights are provided by \citeauthor{penkova_2016} (\citeyear{penkova_2016}; \citeyear{Pen’kova2017}; \citeyear{penkova_2018}; \citeyear{Pen’kova2019b}; \citeyear{Pen’kova2019a}; \citeyear{Pen’kova2020}).}




\subsubsubsection{South Slavic}\label{sec:wiemer:3.2.1.1}

In South Slavic, its middle Section, the Štokavian-Čakavian area, is most interesting. For this area, the most comprehensive account to date is \citet{Grković-Mejdžor2012}. She reports (\citeyear{Grković-Mejdžor2012}: 95--97) that the original construction \textit{bud}- + \textit{l}-form was extremely rare in its western (i.e. the Čakavian) part,\footnote{\label{fn:22}However, according to an anonymous reviewer, \citet[350]{Lukežić2015} cites \textit{bud}- + \textit{l}-participle (alongside with the infinitive!) for sources dating back to the 13th century. I am unable to check this (see fn. \ref{fn:18}).} while in the eastern (Štokavian) part it was attested in the 14\textsuperscript{th} century. It could still have an anterior function, like a future perfect, but it also occurred in hypothetical conditionals and in contexts with past reference. Thus, the anterior function ceased to be the core function, but it might have been the functional pivot which explains why, in some of the early documents, \textit{bud}- + \textit{l}-form is also occasionally attested as an epistemic downtoner or reportive marker. Otherwise, the construction simply turned into a future. 

In the western part, instead of the \textit{l}-form, we find \textit{bud}- in combination with the infinitive from the end of the 14th century (see also fn. \ref{fn:22}). \citet[99]{Grković-Mejdžor2012} calls it a Štokavian-Čakavian innovation. The construction first appeared in dependent clauses, where, in the majority of cases, it marked precedence to another situation. In independent clauses \textit{bud}- + infinitive first appeared in Čakavian (from the second half of the 14\textsuperscript{th} century), before it spread into the remaining (Štokavian) territory, i.e. from west to east within that area, but not beyond it: over the entire history of Balkan Slavic \textit{bǫd}-/\textit{bud}- + infinitive is unattested (\cite{Ivanova-Mirčeva1962}: 144).

Consequently, in early Čakavian and Štokavian varieties (14\textsuperscript{th}--15\textsuperscript{th} centuries) we observe a somewhat unclear diatopic distribution of \textit{bud}-+ infinitive (innovative) and \textit{bud}- + \textit{l}-form (inherited). Nonetheless, at least Grković-Medjžor’s account suggests that the innovative pattern from the west and the older pattern from the east “ran into one another”, causing different “mixtures” in various pockets, as this shows up today in some dialects (see above). In both patterns, stems of either aspect occurred, although for the new pattern (\textit{bud}- + infinitive) ipfv. stems were 2--3 times more frequent than pfv. ones (\cite{Grković-Mejdžor2012}: 97). This was practically inverse to the proportions of pfv. vs ipfv. stems in the original pattern \textit{bud}- + \textit{l}-form (see below). It therefore seems that in South Slavic for either pattern there never was any restriction as for aspect. However, for various historical and contemporary local varieties it needs to be checked to what extent PFV.PRS was/is likewise widespread to mark future reference (also in main clauses) and, thus, whether the distribution of \textit{bud-/bô}- was/is less symmetric for the aspect of the lexical verb on account of PFV.PRS being used as an alternative to \textit{bud-/bô}- with pfv. stems.

We are left with the question how the original pattern \textit{bud}- + \textit{l}-form “survived” in Slovene, i.e. the northwest corner of South Slavic. Slovene remained unaffected by the “Balkanized” type of future, which, based on successors of want-verbs, was an innovation from the southeast (see Section \ref{sec:wiemer:3.2.2}), but this does not explain why in Slovene the pattern \textit{bô}- + \textit{l}-form has persisted as the only one. Unfortunately, the earliest Slovene texts (\textit{Freisinger Denkmäler}, 10\textsuperscript{th}--11\textsuperscript{th} centuries) show no evidence of *\textit{bǫd}- at all (\cite{Whaley2000}: 27--28), nor are there any attestations of *\textit{bǫd}- + infinitive in the history of Slovene (\cite{andersen_2006b}: 19). 



\subsubsubsection{North Slavic}\label{sec:wiemer:3.2.1.2}

Concerning the development of the `become'-based future in West Slavic, the most comprehensive and thorough account, based on a large comparative body of evidence, is still \citet{Křížková1960}, regardless of criticism (see Section \ref{sec:wiemer:3.2.1.3}). Within North Slavic, \textit{bǫd}-/\textit{bud}- most probably started to consolidate as a future marker in Bohemia (13\textsuperscript{th} century). \citet[92 et passim]{Křížková1960} emphasizes that \textit{bud}- + infinitive figures in Old Czech texts from the end of the 13th c. as a largely established construction, which outnumbered constructions with \textit{jmati} ‘have’ and \textit{chtíti} ‘want’ (+ infinitive of either aspect). The latter occurred much more rarely in futurate contexts and rather retained their modal or volitional functions. The pattern \textit{bud}- + infinitive also superseded verbs with ingressive prefixes, which in the earliest attested layers of Czech served to mark “durative” future (see Section \ref{sec:wiemer:4}). Compare, e.g., \textit{ktož \textbf{vz-}bydlí}.\textsc{prs.3sg}\textit{ v stanu tvém }‘who \textbf{starts to live} (> will live) in your camp‘, later “replaced” by \textbf{\textit{bud-e}}.\textsc{fut.3sg}\textit{ bydli-ti}.\textsc{ipfv.inf} (\cite{Křížková1960}: 92--94). 

Further Křížková reports that the Old Czech texts provide almost no evidence of a \textit{bud}- + \textit{l}-form, and only in some of the few attested cases could this pattern be interpreted as a future perfect.\footnote{A slightly larger number of \textit{bud}- + \textit{l}-form can be found only in the \textit{Knihy rožmberské} (1360), but even in these legal texts there are 8 attestations with \textit{l}-form (against 11 with infinitive), from which 7 are from pfv. stems. This corresponds to a higher number of cases of\textit{ bud- + l}-participle in early East Slavic law texts (e.g., the \textit{rus’kaja pravda}) (see below).} In contrast to East Slavic (see below), no modal functions were attested (\cite{penkova_2016}: 481). Moreover, the small empirical basis does not allow us to build any plausible link to \textit{bud}- + infinitive, all the more as the few \textit{l}-forms predominantly built on pfv. stems, while right from the start infinitives with \textit{bud}- exclusively belonged to ipfv. stems. Obviously, these were two different constructions with complementary choice of aspect in the lexical verb (\citeyear{Křížková1960}: 95--96). Shortly after, the construction with the \textit{l}-form disappeared altogether.

In the Sorbian language area, the construction \textit{bud}- + \textit{l}-form has not been attested at all, while \textit{bud}- + infinitive (with either aspect) started occurring from 1500 (\cite{andersen_2006b}: 16).

As for Old Polish, the situation basically differs only with respect to the fate of the \textit{l}-form. In the earliest attestations (mid-late 14th c.) \textit{będ- + l}-form is rare, the infinitive prevails, but \textit{l}-forms already occur exclusively from ipfv. stems\footnote{Sporadically \textit{l}-forms of pfv. stems appear in the 16\textsuperscript{th} century, but only in translations from Ruthenian texts (e.g., from the \textit{Statuty Litewskie}), e.g. \textit{będzie zakupił }‘he will have bought’,\textit{ będą dali} ‘they will have given’ (\cite{Křížková1960}: 106). These could probably better be interpreted as epistemic judgments (see below).} and not as a future perfect (\cite{Křížková1960}: 106, 156); see \xref{ex:wiemer:17}:

\ea \label{ex:wiemer:17}
    \gll \textit{a} \textit{\textbf{będzi-e}} \textit{\textbf{krolowa-ł}} \textit{w} \textit{domu} \textit{Jakobowym} \textit{na} \textit{wieki}\\
         and \textsc{fut-3sg} reign[\textsc{ipfv}]-\textsc{lf-(sg.m)} in house of.Jacob on ages\\
    \glt ‘And he \textbf{will reign} in the house of Jacob for ages.’\\
    \hfill (Ew\_Zam\_Mat[1\_2r]\_33; from \cite{Błaszczak2019}: 187)
\z

In the mid-15\textsuperscript{th} century, the \textit{l}-form reappears (\cite{Křížková1960}: 106). Most probably, this “renaissance” of the \textit{l}-form in the ipfv. future of Polish does not directly continue the CS construction \textit{bǫd}- + \textit{l}-participle (see example 16 from OCS). In the earliest extant texts, infinitival complements clearly prevail (approx. 90\% of all attestations). Only in the 16\textsuperscript{th} century does the frequency of \textit{l}-forms increase again, but it is secular texts which forward their spread in the periphrastic future. It seems as if the \textit{l}-form in this construction reappeared after a period of oblivion, and for that time the distribution (\textit{l}-form vs infinitive) resembled register variation (\cite{Křížková1960}: 104, \cite{Błaszczak2019}: 208--213). Subsequently this variation was levelled out (see fn. \ref{fn:19}).

Let us turn to East Slavic. In general, within Slavic, this was the area with the strongest propensity for ingressive verbs (in their present tense forms: \textit{učьn}\babelhyphen{nobreak}, \textit{počьn}\babelhyphen{nobreak} and \textit{jm}\babelhyphen{nobreak} ‘begin’, etc.) to serve as quasi-auxiliaries of future (with ipfv. infinitives) until the late 16\textsuperscript{th} century (cf. \cite{Pen’kova2019a} for a corpus-based evaluation). This was in contrast to South Slavic, where such verbs hardly played any role as candidates of future markers (\cite{Křížková1960}: 166, 174--175, \cite{Ivanova-Mirčeva1962}, \cite{Grković-Mejdžor2012}). In turn, \textsc{want}- and \textsc{have}-based periphrases (with infinitives) marking future events in early East Slavic texts have been considered as reflecting Church Slavonic (i.e. eastern South Slavic) influence in written language (\cite{andersen_2006b}: 27; \citeyear{Andersen2006a}: 72--73, 81).

In turn, the construction \textit{bud-} + \textit{l}-form is well-attested in the earliest East Slavic texts (11\textsuperscript{th}--12\textsuperscript{th} century) and can be considered original. The oldest documents show \textit{bud}- + \textit{l}-form with stems of either aspect, although pfv. stems prevail. It was more typical of administrative language. It was not restricted to conditional or other dependent clauses, and only rarely did it function as a “simple” future, but also its function as a future perfect has been disputed. In the Old Novgorod dialect \textit{bud}- + \textit{l}-form lost its function as a relative tense (future perfect) early and tended towards acquiring an inferential function (\cite{Pen’kova2019b}).\footnote{\citet[76--80, 82--84]{Andersen2006a} took issue with the commonplace analysis of relevant examples in the Novgorodian birchbark letters and claimed that it was wrong. However, the bulk of examples are in conditional or temporal clauses, for which one cannot clearly decide whether reference is made to a single event and, above all, from whose perspective temporal relations should be judged (the writer’s, the addressee’s, or of an independent observer?).} In northern Ruthenian\footnote{\textit{Ruthenian} refers to those parts of East Slavic which from the late 14\textsuperscript{th} century came under increasing Polish influence (within the \textit{Rzeczpospolita Obojga Narodów}) and which can be considered the soil on which later Ukrainian and Belarusian were formed.} documents of the 14\textsuperscript{th}--16\textsuperscript{th} century, \textit{bud}- + \textit{l}-form practically got restricted to contexts with suspended propositions, where it acquired dubitative and reportive functions, mainly in dependent clauses (\cite{Pen’kova2020}). Eventually, throughout East Slavic, the construction disappeared in the 16\textsuperscript{th} century, leaving behind the \textsc{3sg}-form\textit{ budetъ} as an epistemic or inferential particle (\cite{Křížková1960}: 162--169, \cite{Pen’kova2019a}); see also Section \ref{sec:wiemer:3.2.1.3}.

The construction \textit{bud}- + infinitive is attested later, and practically only with ipfv. stems. It first appears in the second half of the 14\textsuperscript{th} century (in parallel to \textit{imati} ‘have’ and phasal verbs) in Ruthenian texts, i.e. at the (south)western periphery of the East Slavic territory. In texts of the (north)eastern East Slavic territory (“Muscovy”) it only appears from the mid-15\textsuperscript{th}c., among others (but not exclusively) in written exchange with Polish diplomats and similar official documents. Only during the 17\textsuperscript{th} century does \textit{bud}- + infinitive gain more prominence: through the first half of the 17\textsuperscript{th} century it is still rare (much rarer than \textit{učn-} or \textit{stan-} ‘begin’) in the administrative language, but its use increases in the second half of the 17\textsuperscript{th} century, until it finally ousts \textit{učn-} (but not \textit{stan-}) and becomes representative of daily speech (\cite{Křížková1960}: 173--176). Through the 17\textsuperscript{th} century the spread of \textit{bud}- was clearly supported by a Polish model, only by the 18\textsuperscript{th} century did \textit{bud}- + infinitive consistently appear in texts without Polish influence (\cite{Moser1998}: 311--330). 



\subsubsubsection{Where did the infinitive come from?}\label{sec:wiemer:3.2.1.3}

Let us turn to the question how the infinitive became integrated into a construction with \textit{bud}-. Different theories have been developed, all of which require qualification.

In the first place, how could an originally inchoative verb combine with an infinitive, either before or after the inchoative feature was lost? Concerning Old Serbian (eastern Štokavian), \citet[97]{Grković-Mejdžor2012} suggests that the infinitive introduced an expectative viewpoint, but she does not attempt to explain how the expectative (or goal-oriented) semantics of an infinitive might combine with an originally inchoative verb.

Actually, from a purely semantic point of view, the combination of an inchoative verb (inducing a change of state) with a goal-directed form (like the infinitive) does make sense. However, upon closer examination, diachronic Slavic syntax does not supply any good model for a combination of \textsc{be}/\textsc{become} + infinitive of purpose as an emergent future. Admittedly, constructions with a copula and an infinitive, with the latter contributing a meaning of purpose (due to its IE provenance), are well-attested in Slavic (and Baltic), and so are possessive constructions of the \textsc{mihi est liber}-type, i.e. with the possessor in the dative and the possessee in the nominative. Such constructions laid the basis for modal (deontic or circumstantial) constructions, particularly in Russian (e.g., \textit{Mne}.\textsc{dat} \textit{tut nečego delat’}.\textsc{inf} ≈ ‘There is nothing for me to do here’), but also in historical documents, such as the famous example \xref{ex:wiemer:18} from the \textit{rus’kaja pravda}, the oldest kind of legal codex in East Slavic (11\textsuperscript{th} century):

\ea \label{ex:wiemer:18}
    \gll \textit{takov-a} \textit{pravd-a} \textit{uzja-ti} \textit{rusin-u}\\
         such-\textsc{nom.sg.f} right[\textsc{f}]-\textsc{nom.sg} take[\textsc{pfv}]-\textsc{inf} Rus’.inhabitant-\textsc{dat.sg}\\
    \glt ≈ ‘These are the rights a “Russian” must/can enjoy.’\\
    \hfill (\cite{Holvoet2003}: 467, translation adapted)
\z

However, in these constructions the ``possessor" (i.e. more agent-like argument) is coded in the dative (cf. \cite{Holvoet2003} for an analysis and the diachronic background), whereas in the \textit{bud}- + infinitive-construction the corresponding argument appears in the nominative.\footnote{By comparison, even in languages like German modal infinitive constructions only code the more patient-like argument in the nominative, while the potential agent normally is left out (e.g., \textit{Das Problem ist (nicht) zu lösen}.\textsc{inf} ‘The problem is (not) to be solved’).} This is also admitted by \citet[34]{andersen_2006b}, who points out a construction with \textit{bude} or \textit{budetь} as an existential verb; these forms are encountered in the Novgorod birchbark letters (see \ref{ex:wiemer:19}). Equivalent examples are known from 13\textsuperscript{th} century documents farther to the south (see \ref{ex:wiemer:20}).

\ea \label{ex:wiemer:19}
    \gll \textit{ažь} \textit{dolgo} \textit{\textbf{bude}} \textit{dolgo} \textit{medlę-ti} \textit{prisl-i} \textit{vest-e}\\
         if long \textsc{fut.3sg} long overstay-\textsc{inf} send-\textsc{imp.sg} message-\textsc{acc.sg}\\
    \glt ‘If it \textbf{happens} (for you) to stay much longer, send (me) a message.’\\
    (Gramota no. 771, end of 13\textsuperscript{th}--beginning 14\textsuperscript{th} century; \cite{Zaliznjak2004}: 532)
\z

\ea \label{ex:wiemer:20}
    \gll \textit{aže} \textit{\textbf{budětь}} \textit{tъrgova-ti} \textit{smolnęnin-u} \textit{sъ} \textit{nemьčic-emь}\\
         if \textsc{fut.3sg} trade-\textsc{inf} Smolensk.inhabitant-\textsc{dat.sg} with German-\textsc{ins.sg}\\
    \glt ‘If it \textbf{happens} for an inhabitant of Smolensk to trade with a German.’\\
    \hfill (contract, Smolensk, 1220s; \cite{Zaliznjak2004}: 532)
\z

Despite the gloss \textsc{fut.3sg}, here \textit{bude/budětь} is to be considered an invariable form employed either like a clausal operator or as a predicate which takes a clausal complement with an infinitival head. Again, if mentioned at all, the most agent-like argument of the infinitive occurs in the dative (see \ref{ex:wiemer:20}). Such examples clearly show the use of petrified \textit{bude/budětь }in habitual or conditional contexts, and it certainly was this use which paved the way for \textit{budetъ} as an epistemic-inferential particle after the 16\textsuperscript{th} century (see Section \ref{sec:wiemer:3.2.1.2}). As for the structure, \citet[32, 35]{andersen_2006b} only surmises that it might have arisen as ``a syntactic neologism coined for subjectless verbs" (as in some hypothetical CS *\textit{bǫdetъ moroziti }‘(it) will freeze’), which later expanded to clauses with nominative subjects. However, this is only an \textit{ad hoc}-stipulation, so that Andersen ultimately concedes that there is no real solution to the issue.

In sum, neither in East Slavic, nor anywhere else in Slavic, are any combinations of \textit{bud}- + infinitive in purpose meaning attested even in the earliest documents. Apart from this, there is no reason to expect a free combination of \textit{bud}- with infinitive to restrict itself to, or prefer, ipfv. aspect of the infinitive. However, from the earliest documents onwards, only ipfv. infinitives are attested in West Slavic (Section \ref{sec:wiemer:3.2.1.2}) and ipfv. infinitives considerably outnumbered pfv. ones in Old Serbian (Section \ref{sec:wiemer:3.2.1.1}).

Alternatively, \citet[98]{Křížková1960} argues that \textit{bud}- + infinitive acquired future meaning holistically, as a construction (with loss of internal transparency). However, since she also rightly discards any connection to possessive constructions (see above), one wonders how this holistic meaning and the implied loss of inchoativity of\textit{ bud}- came about. At any rate, these processes must have taken place prior to the earliest documents in Old Czech (i.e. during the early 13\textsuperscript{th} century or even earlier).

Finally, an entirely different approach was chosen by \citet{blaszczak_etal_2014}. The authors decompose the construction into a pfv., possibly ingressive (‘inceptive’) operator (\textit{będ}-) and its ipfv. ``complement'' (infinitive or \textit{l}-form). Since the pfv. operator is in the present (or non-past) tense, it causes a “forward shift” of the reference time, which the authors take to be equivalent to a shift to future (\citeyear{blaszczak_etal_2014}: 190--191). They further claim that \textit{będ}- denotes a Kimian state. Such a state does not represent an eventuality, but ``a referential argument for a temporally bound property exemplification" (\citeyear{blaszczak_etal_2014}: 190). It is located in time but not in space and exists only for the needs of language communication, as in \textit{Bill knew the answer (and that’s why he passed the exam)} (\cite{Tatevosov2016}: 63--65). \citeauthor{blaszczak_etal_2014} do not justify why \textit{będ}- should be assigned this status, nor do they notice the inherent contradiction when they treat \textit{będ}- as the pfv. head of the ASP-slot in syntactic structure and simultaneously claim it to represent a state. One might assume that \textit{będ}- marks the beginning of a state (which would be compatible with the assumed inchoative origin of *\textit{bǫd}-), but this obviously is not what the authors accept, either, as it would not conform to their own examples. These examples, in turn, mostly exemplify processes (not states), which are however not distinguished by these authors. Finally, and after all, \citeauthor{blaszczak_etal_2014} only concentrate on the question why \textit{będ}- exclusively combines with ipfv. stems of the lexical verb, but they do not clarify how the infinitive “entered the scene”, let alone why in all North Slavic languages except Polish and Kashubian it entirely ousted, or replaced, the \textit{l}-form. 

To my knowledge, there have been no more attempts at explaining the rise of \textit{bud-/będ}- + infinitive on internal grounds. However, since the 1950s there have been hot debates concerning the issue whether the \textit{bud}-future construction might have been PAT-borrowed (in the sense of \cite{MatrasSakel2007}) from Middle or Early New High German into Old Czech, or the other way around. Remarkably, this discussion mostly disregards this construction in South Slavic (see Section \ref{sec:wiemer:3.2.1.1}), and hardly anybody thought about the possibility of mutual influence between western edges of Slavic and some German varieties (also closer to Slovene and CentSS). As for German, the earliest attestations of \textit{werdan} + infinitive have usually been dated between the early 11\textsuperscript{th} century and the early 13\textsuperscript{th} century (\cite{Westvik2000}). However, this construction was also attested with \textit{werdan} in the past tense in the 13\textsuperscript{th}-16\textsuperscript{th} centuries (\cite{Kraemer2005}: 75--81). Jointly, \textit{werden} + infinitive seems to have kept an inchoative meaning until the time of its alleged influence on *\textit{bǫd}- + infinitive in West Slavic. However, *\textit{bǫd}- + infinitive could never occur in the past tense, and *\textit{bǫd}- itself must have lost its inchoative flavour in Czech no later than during the 13\textsuperscript{th} century.\footnote{For some more details and a concise evaluation of these debates cf. \citet[105--108]{wiemer_hansen_2012}. For surveys on the German \textit{werden}-future cf. \citet[Section \ref{sec:wiemer:2}]{Hartmann2021} and \citet[Section \ref{sec:wiemer:2}]{Jäger2024}.}

Recently, \citet{Jäger2024} has suggested that the genesis of the German \textit{werden}-future started with the ``reanalysis of originally predicative past participles in \textit{werden}-copula constructions as infinitives''. This reanalysis was possible because the past participles of a small class of strong unaccusative verbs lacked a \textit{ge}-prefix. These forms were homonymous with their infinitives (compare modern Germ. \textit{Das Gebäude wird\textbf{ zerfallen}}.\textsc{ptcp} \textit{sein} ‘The building will (or: may) be fallen apart’ vs \textit{Das Gebäude wird \textbf{zerfallen}}.\textsc{inf} ‘The building will fall apart’), so that as predicative “complements” of inchoative \textit{werdan} ‘become’ in the present tense these participles created bridging contexts for the reanalysis of a future-oriented resultative construction as a future. According to Jäger, the reanalysis occurred in Old High German and Old Low German, i.e. until the 11\textsuperscript{th} century. Thus, if her hypothesis turns out to be correct, the reanalysis should have been followed by analogical expansion, before the new \textsc{become}+infinitive construction could have served as a model for Slavic varieties at the western edges. That is, this construction should have already spread considerably to become sufficiently established in Middle High German prior to being copied by speakers of Slavic, and one wonders why speakers of Old Slovene did not adopt this construction. 

In sum, the issue whether PAT-borrowing or, alternatively, polysemy copying\footnote{For a discussion of these notions cf. \citet[26--29]{wiemer_walchli_2012}.} played a role as a catalyst of the spread of \textit{bǫd-/bud}- + infinitive in Slavic is still open. 




\subsubsection{\textsc{want}}\label{sec:wiemer:3.2.2}

CS *\textit{xotěti/xătěti} ‘want’ is the predominant source of future marking in contemporary South Slavic, but its embryonic stages can be observed also in North Slavic, where, apart from marking intention and wish, it can at best be regarded as a prospective marker (\cite{Křížková1960}: 120--125, \cite{seveleva2021}). Remarkably, it does not seem to have ever shown any preference for the aspect of the stem of the lexical verb. In this respect, modern South Slavic languages (except Slovene) probably just continue the distribution from the onset of the grammaticalization of \textsc{want}-periphrases.

In early OCS, \textsc{want}-periphrases referring to future events have been claimed to be initially less frequent than equivalent \textsc{have}-periphrases (\cite{Birnbaum1958}, \cite{Grković-Mejdžor1995}, also \cite{Xaralampiev1981}: 116 with further references). In the eastern part of South Slavic, this relation must have swapped slightly later, after or even during the 11\textsuperscript{th} century. As for the Štokavian area, \citet{Grković-Mejdžor2012} concludes that from the mid-13\textsuperscript{th} century the present tense forms of \textit{xtěti} ‘want’ started showing signs of phonetic erosion, cliticization (> \textsc{2p}-enclitic) and context expansion to inanimate situations (with concomitant semantic bleaching). This process had already basically been finished in the 14\textsuperscript{th} century; it seems to have lagged behind slightly in the Čakavian area, i.e. further to the Adriatic coast in the west (\cite{Grković-Mejdžor2012}: 88--89). 

As concerns early Slavic in the Balkans, \textit{xotěti} must have already started taking over the floor during the 11\textsuperscript{th} century (\cite{tomic_2004}: 409--410, referring to \cite{Asenova2002}: 204 and \cite{Mirčev1978}: 223). \citet[191]{HeineKuteva2005} report that \textsc{want} competed with \textsc{have} and \textsc{begin} until the 13\textsuperscript{th} century, when \textit{xošte} (> \textit{šta}) + truncated infinitive gained priority. From the 15\textsuperscript{th} century onwards the infinitive was gradually replaced by \textit{da}-complements (with finite verbs), which then also appeared after \textit{šta > šte} (\cite{tomic_2004}: 405, referring to \cite{Koneski1967}: 206 and \cite{Asenova2002}: 205). \citet[407--408]{tomic_2004} argues that the \textit{da}-construction bridged the transition from an inflected “modal clitic” (\textsc{want}) with infinitive towards the contemporary construction (uninflected proclitic + V\textsubscript{\textsc{prs}}). Table \ref{tab:wiemer:1} (based on \cite{Asenova2002}: 214) summarizes these stages in a simplified manner with the ipfv. stem \textit{pisa-/piš}- ‘write’ and its \textsc{2sg}-form (of the present tense stem \textit{piš}-) for the purpose of illustration.

\newcolumntype{L}{>{\raggedright\arraybackslash}X}

\begin{table}
\footnotesize
\caption{From \textsc{want} to future in South Slavic (here for Bulgarian) }
\label{tab:wiemer:1}
 \begin{tabularx}{\textwidth}{LLLLL}
   \lsptoprule
9th--13th c. (OCS) & 14th--20th c. & 15th--early 20th c. & 16th--21st c. & 16th--21st c.\\
  \midrule
\textit{xošt-eši pisa-ti} & \textit{xošt-eš pisa} & \textit{xošt-eš da piš-eš} & \textit{šta \textgreater{} šte da piš-eš} & \textit{šte piš-eš}\\
\midrule
want-\textsc{prs.2sg} + \textsc{inf} & want-\textsc{prs.2sg} + truncated \textsc{inf} & want/\textsc{irr}-\textsc{prs.2sg} + \textsc{irr} + \textsc{prs.2sg} & \textsc{fut}/\textsc{irr} (uninfl.) + \textsc{irr} +\textsc{prs.2sg} & \textsc{fut} (uninfl.) + \textsc{prs.2sg}\\ 
  \lspbottomrule
 \end{tabularx}
\end{table}

What is remarkable about this sequence of stages is not so much the gradual phonological and categorial deterioration of the initial \textsc{want}-verb (\textit{xošte-ši}.\textsc{prs.2sg} > \textit{šta} > \textit{šte}), nor the loss of the infinitive (``replaced" by inflected present tense forms) in the lexical verb, but the suggested immense overlap among the future marker variants and their combinations with and without the irrealis proclitic \textit{da}.\footnote{Cf. \citet[134 and passim]{Ivanova-Mirčeva1962} for details. She points out that truncated forms (\textit{štǫ}.\textsc{1sg}, \textit{šteši}.\textsc{2sg}, etc.) could be found as early as in the 13th century, while contemporary uninflected \textit{šte} became predominant only by the 18th century (\citeyear{Ivanova-Mirčeva1962}: 134--135).} Here, again, one has to be careful about the data used: most probably, these overlaps rather reflect an enormous diatopic and diaphasic variation for most, if not all, centuries; this analysis ``overrides" rather innovative tendencies in vernacular speech with conservative text genres based on OCS writing traditions. The modern standard form of the future without \textit{da} (\textit{šte pišeš}) gained dominance only quite recently, whereas the pattern with \textit{da} (\textit{šte da }\textit{pišeš}) has been marginalized as a device of marking assumptions (\cite{asenova_guentcheva_2022}).

Regardless, \textit{da}-V\textsubscript{\textsc{prs}} has persisted in the suppletive \textsc{have}-pattern which dominates in the negated future (\textit{njama da} + V\textsubscript{\textsc{prs}}). In this construction, \textsc{have} occurs in the invariable form of \textsc{prs.3sg} (see \ref{ex:wiemer:21b} and \ref{ex:wiemer:22}). In Macedonian we encounter an analogical situation\footnote{\citet[388]{Tomić2012} remarks that the construction with negated \textit{ḱe }conveys an additional “presuppositional component” (compare, e.g., \textbf{\textit{Ne} \textit{ḱe}}\textit{ te zemat so niv }‘(I presume) they \textbf{won’t} take you with them’ with \ref{ex:wiemer:21b}).} (see \ref{ex:wiemer:9}). 

\ea Macedonian \label{ex:wiemer:21}\\
    \ea \label{ex:wiemer:21a}
    \gll \textit{Nie} \textit{ḱe=} \textit{stign-eme} \textit{utre.}\\
         1\textsc{pl.nom} \textsc{fut=} arrive[\textsc{pfv}]-\textsc{prs.1pl} tomorrow\\
    \glt ‘We will arrive tomorrow.’
    \ex \label{ex:wiemer:21b}
    \gll \textit{Nema} \textit{da=} \textit{te=} \textit{zema-m.}\\
         \textsc{neg}.have.\textsc{invar} \textsc{irr=} 2\textsc{sg.acc} take[\textsc{pfv}]-\textsc{prs.1sg}\\
    \glt ‘I won’t take you (with me).’
    \z
\z

In modern standard Macedonian and Bulgarian we also find the unnegated pattern \textsc{have} + \textit{da} + V\textsubscript{\textsc{prs}}, which still carries some obligational flavour; compare \xref{ex:wiemer:10}:

\ea Macedonian \label{ex:wiemer:22}\\
    \gll \textit{Ima} \textit{da=} \textit{oda-m.}\\
         have.\textsc{invar} \textsc{irr=} go[\textsc{ipfv}]-\textsc{prs.1sg}\\
    \glt ‘I have to go > I will go.’
\z

These patterns suggest a persistent “rivalry” between \textsc{want}- and \textsc{have}-based constructions showing relation to future marking, for which the degree of entrenchment varies. However, their exact diachronic relationship, in particular of the different \textsc{have}-based future periphrases in the Balkans (is it a direct continuation or a cyclic phenomenon?) requires clarification (cf. \cite{Xaralampiev1981}: 119--122, \cite{kramer1997negation}: 408, \cite{mitkovska_buzarovska_2014}: 195--196). As for Old Serbian, the analysis presented by \citet[93--97]{Grković-Mejdžor2012} even seems to indicate that here the history of predicative \textsc{have}-constructions and their emergent development as possible futures started anew only in the 14\textsuperscript{th} century, which would mean that it was entirely detached from \textsc{have}-based future-like constructions in OCS.

Most of the aforementioned processes, and the variation, have parallels in other languages of the Balkans. We observe them in the history of the Greek \textsc{want}-future, and the diachronic competition between \textsc{want} and \textsc{have} finds their diatopic equivalent in contemporary differences between North (Geg) and South (Tosk) Albanian. For overviews cf. \citet{kramer1994grammaticalization}, \citet[36--38]{andersen_2006b}, \citet{Tomić2006}, and \citet[51--53]{WiemerLinnemeier2024}.

\section{Future marked by prefixes on imperfective stems}\label{sec:wiemer:4}

In Czech (\cite{Šmilauer1947}: 152, \cite{Machek1962}: 439, \cite{Mrhačová1993}: 65--67, among many others), Slovak (\cite{Horecký1957}, \cite{RužičkaDvonč1966}: 41), in Upper and Lower Sorbian (\cite{Michałk1963}, \cite{Scholze2008}: 194--195, 269--270, 279--280) and (to a lesser extent) in Slovene we observe a peculiar phenomenon. The prefix \textit{po}- (in Sorbian also \textit{s-/z-}) marks future tense of certain classes of ipfv. stems. The core group has always consisted of verbs of directed motion (e.g., Cz. \textit{pojedu} ‘I will go/drive’, \textit{popluji} ‘I will swim’), related causatives (e.g., \textit{potáhnu} ‘I will pull’, \textit{povezu} ‘I will bring (by driving)’), and inchoative verbs (e.g., \textit{poroste} ‘s/he will grow’), occasionally also some non-motion activity verbs (e.g.,\textit{ mlýn pomele} ‘the mill will grind’, \cite{Kopečný1962}: 49) and even state verbs (e.g., colloquial USorb. \textit{změjem} ‘I will have’, \cite{Scholze2008}: 194) are mentioned. The extension of these groups has remained fuzzy; diachronically it has most probably been subject to variation (cf. \cite{Šlosar1981}: 79--80, 102--103 on Czech). In the modern languages we cannot assume closed classes either. Figures adduced, e.g., for contemporary Czech vary from 20 to 115 lexical items (\cite{Kopečný1962}: 49--50, \cite{Osolsobě2014}: 137, \cite{ŠevčíkováPanevová2018}: 179--180, \cite{VAGSČ2021}: 156, with further references). Before dealing with diachronic issues (Section \ref{sec:wiemer:4.2}--\ref{sec:wiemer:4.3}), I will survey the contemporary stage.


\subsection{Contemporary stage}\label{sec:wiemer:4.1}

First of all, ipfv. stems with \textit{po}- marking future typically denote parametric changes, i.e. changes along scales without absolute endpoints (contrary to strictly telic verbs, which imply an absolute goal). Moreover, we are dealing with stems in particular meanings here. First, the ipfv. behaviour of \textit{po}-prefixed future forms often does not apply to these verbs in non-physical, figurative usage, although this is not a strict rule (see below). Second, \textit{po}-prefixation with the same ipfv. stems turns these stems into pfv. aspect if the number of arguments increases (thus, from intransitives to transitives and from transitives to ditransitives). For instance, \textit{pojít} in the meaning ‘originate, rise’ or ‘perish’, or \textit{povléct} in the meaning ‘cover (some surface with sth)’ have a full verbal paradigm and behave like other pfv. stems;\footnote{Cf. \citet[544--545]{Bondarko1961}, \citet[238, 248]{Němec2009a} and most systematically \citet[181--186]{ŠevčíkováPanevová2018}.} compare (\ref{ex:wiemer:23}--\ref{ex:wiemer:24}):

\ea \label{ex:wiemer:23} Czech\\
    \ea \label{ex:wiemer:23a}
    \textit{\textbf{Půjde}}\textsuperscript{\textsc{ipfv.fut}} / \textit{\textbf{šel}}\textsuperscript{\textsc{ipfv.pst}} \textit{domov} \textit{celou} \textit{hodinu.}\\
    ‘He \textbf{will go / went} home for a whole hour.’
    \ex \label{ex:wiemer:23b}
    \textit{Myš} \textit{\textbf{pojde}}\textsuperscript{\textsc{pfv.fut}} / \textit{\textbf{pošla}}\textsuperscript{\textsc{pfv.pst}} \textit{po} \textit{požití} \textit{jedu} \textit{*celou} \textit{hodinu.}\\
    ‘After the mouse eats / ate the poison, it will die / died *for a whole hour.’\\
    \hfill (from \cite{ŠevčíkováPanevová2018}: 182, with adaptions)
    \z
\z

\ea \label{ex:wiemer:24} Czech\\
    \ea \label{ex:wiemer:24a}
    \textit{Kůň} \textit{\textbf{povleče}}\textsuperscript{\textsc{ipfv.fut}} / \textit{\textbf{vlekl}}\textsuperscript{\textsc{ipfv.pst}} \textit{kaskadéra} \textit{za} \textit{sebou} \textit{dlouhé} \textit{minuty.}\\
    ‘The horse \textbf{will drag / dragged} the stuntman after itself for many minutes.’
    \ex \label{ex:wiemer:24b}
    \textit{\textbf{Povleče}}\textsuperscript{\textsc{pfv.fut}} / \textit{\textbf{povlékla}}\textsuperscript{\textsc{pfv.pst}} \textit{postel} (\textit{čerstvým} \textit{prádlem}) \textit{*dlouhé} \textit{minuty.}\\
    ‘She \textbf{will cover / covered} the bed (with fresh linen) *for many minutes.’\\
    \hfill (from \cite{ŠevčíkováPanevová2018}: 182, with adaptions)
    \z
\z


Basically, the same applies to other languages of the “western hemisphere” of Slavic, e.g. Slovene (\cite{dickey_2018}: 375). 

In modern Czech, this kind of lexical split causes grammatical homonymy of \textit{po}-prefixed forms, as can be demonstrated with the following examples (in \ref{ex:wiemer:25a}) and their past equivalents (in \ref{ex:wiemer:25b} taken from the Czech National Corpus (\href{https://www.korpus.cz/}{https://www.korpus.cz/}, subcorpus SYNv5):

\ea \label{ex:wiemer:25} intransitive, parametric change\\
    \ea \label{ex:wiemer:25a} 
    \textit{Ještě} \textit{dva} \textit{týdny} \textit{se} \textit{Tour} \textit{\textbf{povalila}}\textsuperscript{\textsc{ipfv.fut}} \textit{napříč} \textit{Francií.}\\
    ‘For two more weeks, the Tour \textbf{will roll} across France.’
    \ex \label{ex:wiemer:25b}
    \textit{Dva} \textit{týdny} \textit{se} \textit{Tour} \textit{\textbf{valila}}\textsuperscript{\textsc{ipfv.pst}} \textit{\textbf{(*povalí)}} \textit{napříč} \textit{Francií.}\\
    ‘For two weeks, the Tour \textbf{rolled} across France.’\\
    \hfill (from \cite{ŠevčíkováPanevová2018}: 184)
    \z
\z

\ea \label{ex:wiemer:26} transitive, telic\\
    \ea \label{ex:wiemer:26a} 
    \textit{Prasata} \textit{plot} \textit{\textbf{povalí}}\textsuperscript{\textsc{pfv.fut }}\footnote{Of course, like non-past forms of pfv. stems in general, \textit{povalí} might also be used in dispositional or circumstantial meanings (see Section \ref{sec:wiemer:3.1.2}, \ref{sec:wiemer:3.1.4} and below).} \textit{nebo} \textit{se} \textit{podhrabou.}\\
    ‘Pigs \textbf{will knock down} the fence or dig under it.’
    \ex \label{ex:wiemer:26b} transitive, telic\\
    \textit{Prasata} \textit{plot} \textit{\textbf{povalila}}\textsuperscript{\textsc{pfv.pst}} \textit{(*dlouhý} \textit{okamžik.)}\\
    ‘Pigs \textbf{knocked down} the fence (*for a long moment).’\\
    \hfill (from \cite{ŠevčíkováPanevová2018}: 184)
    \z
\z

Note that, in the given meanings, the ipfv. stems (in \ref{ex:wiemer:23a}, \ref{ex:wiemer:24a}, \ref{ex:wiemer:25a}--\ref{ex:wiemer:25b}) do not have pfv. counterparts, i.e. *\textit{pošel} ‘went’, *\textit{povlékl} ‘dragged’, *\textit{povalil }‘rolled’ do not exist.\footnote{By contrast, the pfv. “homonyms” usually have ipfv. counterparts, e.g. \textit{povléct}\textsuperscript{pfv}\textit{ – povlékat}\textsuperscript{ipfv} ‘cover (surface with sth)’, \textit{povalit}\textsuperscript{pfv}\textit{ – valit}\textsuperscript{ipfv} ‘knock down’; however, ipfv. \textit{pocházet} ‘die, perish’ for pfv. \textit{pojít} is considered obsolete.} They are thus \textit{imperfectiva tantum }lexemes.

This brings us to the question why, in the first place, a prefix attached to the present tense of simplex stems should mark the future of ipfv. aspect, and not the present tense of pfv. aspect (which then, via hypoanalysis, yields future). As for contemporary Czech, \citet[47--50]{Kopečný1962} provides the following arguments:

\begin{enumerate}[label=(\alph*)]
    \item \label{arg:a} The \textit{po}-prefix only occurs with the present tense form, not with the past tense, the infinitive or any other paradigmatic forms of the stem (except the imperative; see below).
    \item \label{arg:b} The forms with \textit{po}- are not used in the narrative present (unlike present tense forms of pfv. stems), nor do they usually occur in modal functions. What \citeauthor{Kopečný1962} (and others) mean by ‘modal’ is the use of non-past forms to mark dispositional or circumstantial meanings; these are regularly associated with meanings that in Slavic linguistics have been referred to under the label of ‘inactual’ or ‘gnomic’ present (see Section \ref{sec:wiemer:3.1.2}). Thus, \citet{Nerad1940} finds that \textit{po}-prefixed ipfv. verbs with future meaning never have this modal meaning. Other scholars, like \citet[529]{Bondarko1961}, \citet[47]{Kopečný1962} and \citet[439--441]{Machek1962}, are less categorical in this regard (see below).
    \item \label{arg:c} The \textit{po}-prefixed forms can also be used for future actions with unlimited duration; thus, not only, e.g., \textit{ty to \textbf{poneseš} dvě hodiny }‘you(\textsc{sg}) (will) \textbf{carry} this for two hours’, but also \textit{ty to \textbf{poneseš} stále }‘you(\textsc{sg}) (will) \textbf{carry} this permanently’ (cf. also \cite{Galton1976}: 45); see also (\ref{ex:wiemer:23a}, \ref{ex:wiemer:24a}). \citet[181]{ŠevčíkováPanevová2018} accept compatibility with relevant adverbials as one of two reliable criteria allowing to determine the aspect (see \ref{ex:wiemer:23}--\ref{ex:wiemer:26}), the other criterion being the synonymy of \textit{po}-prefixed present tense forms with unprefixed past tense forms of the same stem (see \ref{ex:wiemer:25a}--\ref{ex:wiemer:25b} against \ref{ex:wiemer:26a}--\ref{ex:wiemer:26b}).
\end{enumerate}

As for \ref{arg:a}, \textit{po}- is also widely attested with the imperative of the same ipfv. stems, but here \textit{po}- bears the implication that the action is to be performed towards and/or jointly with the speaker (e.g., Cz. \textit{pojď!} ‘come (with me)!’, \textit{polez!} ‘climb (here, to my place)!’).\footnote{\label{fn:35}Imperatives from the corresponding unprefixed stems are neutral in this respect (\cite{Bondarko1961}: 541--542). This suggests that, with the imperative, \textit{po}- has retained some spatial, or deictic, features that predate, or interfere with, the early stages of the rise of the stem\hyp{}derivational aspect opposition (as suggested by \cite{Machek1962}: 438--439).} \citet[438--439]{Machek1962} argues that the prefix \textit{po}- in non-past and in imperative forms derives from different etyma; as indirect proof he points out presumably different cognates in Baltic and earlier Indo-European.\footnote{Compare \textit{par}- ‘home’ (e.g., Lith. \textit{par-grįžti} ‘return home’), from *\textit{pār}, vs \textit{pa}- ‘approaching from below’ (e.g., Lith. \textit{pa-si-rašyti} ‘sign (a paper)’, compare Germ. \textit{unter-schreiben}) from *\textit{upo}. Thus, \textit{pa}- in imperatives would derive from *\textit{pār}, while in non-past forms it would continue *\textit{upo}.}

As for \ref{arg:b}, \citet[439--441]{Machek1962} argues that \textit{po}-, before it came to mark the future of ipfv. stems, might have been used to mark dispositions, or inclinations, of the subject argument (see Section \ref{sec:wiemer:4.3}). In his opinion, this is still reflected in modern Czech; see (\ref{ex:wiemer:27}--\ref{ex:wiemer:28}), which he interprets as remnants of an older layer. Note that \textit{podělat} in \xref{ex:wiemer:28} is not a motion verb, nor does \textit{ponesou} in \xref{ex:wiemer:27} clearly denote directed motion:


\ea \label{ex:wiemer:27} 
    \textit{Do práce budu chodit}\textsuperscript{\textsc{ipfv}} \textit{pěšky tak dlouho, dokud mne nohy \textbf{ponesou}}\textsuperscript{\textsc{ipfv}}.\\
    ‘I will walk to work as long as my legs \textbf{can / will carry} me.’
\z

\ea \label{ex:wiemer:28} 
    \textit{Sám toho \textbf{nepodělám}}\textsuperscript{\textsc{pfv}}.\\
    ‘\textbf{I can’t / won’t do} it alone.’
\z

\citeauthor{Machek1962} points out parallels in Baltic; compare, for instance, Lith. \textit{jis net \textbf{ne-pa-eina}} ‘he \textbf{won’t }even \textbf{go} / he \textbf{is unable to go} (even a bit)’ or \textit{ji \textbf{pa-kelia} bet kokį sunkų lagaminą} ‘she \textbf{is able to lift }any heavy suitcase’.\footnote{According to \citet[441]{Machek1962}, this \textit{pa}- derives from yet another etymon, namely IE *\textit{per}- (> Baltic \textit{par}-).} \citeauthor{Machek1962}'s observation is important, since a diachronic link between these modal (or “timeless”) functions and the future function is very likely, more precisely: the future function might have emerged via abduction from the knowledge about lawful and, thus, predictable behaviour of individuals or of situations (see Section \ref{sec:wiemer:3.1.2}).

However, in today’s Czech, the dispositional function of \textit{po}-prefixed ipfv. verbs is rarely noted. It now appears to be rather associated with a few verb stems prefixed with\textit{ u}- (e.g., \textit{Karel \textbf{unese} 200 kg} ‘Karel \textbf{is able to lift} 200 kg’; cf. \cite{Mrhačová1993}: 62), but these stems are commonly classified as pfv., since they are incompatible with unlimited duration.\footnote{Lexicographic practice seems to treat them as \textit{perfectiva tantum}. Thus, for instance, \textit{unést} does have an ipfv. counterpart \textit{unášet}, but only in the meaning ‘carry away’, not in the slightly different meaning ‘lift’ with a dispositional ``touch''. I am grateful to Jana Kocková for consultation on these issues.}

\largerpage[1.5]

In addition to \citeauthor{Kopečný1962}’s three aforementioned arguments supporting the claim that we are dealing with the future of ipfv. stems, we have to take into account that most of the relevant ipfv. stems can also form their future with \textit{bud}-. However, these stems show this ``periphrastic'' way of expressing ipfv. future (much) more rarely, and they vary considerably among each other as for the likelihood with which we may observe the periphrastic future. Some few stems, like \textit{jít} ‘go’, are not attested with \textit{bud}- at all, while others show larger figures; however the bias is always clearly in favour of prefixation with \textit{po}- (cf., e.g., \cite{Kopečný1962}: 47--48, \cite{VAGSČ2021}: 156, \cite{SaicováŘímalová2023}). It is an open question whether this dual way of marking ipfv. future tends toward functional differentiation. It has been claimed that \textit{bud}- is more likely if the motion verb is used figuratively (e.g., \textit{já mu \textbf{budu lézt} do zadku?!} ‘Me \textbf{crawl} up his ass?!’, cited by \cite[45]{Galton1976} with reference to \cite{Poldauf1948}), but there are also figurative uses with \textit{po}-prefix (e.g. \textit{Tato ulice \textbf{ponese} jméno XY} ‘This street \textbf{will bear} the name XY’, \cite{Galton1976}: 45). Native speakers’ judgments in this regard diverge considerably. Likewise, the analysis of corpus data offered by \citet[181]{ŠevčíkováPanevová2018} and by \citet[]{SaicováŘímalová2023} does not indicate any clear-cut distribution of \textit{po}- vs \textit{bud}- along the lines of, e.g., a differentiation of physical (directed) motion vs figurative use. \citet[]{SaicováŘímalová2023} concludes that, at least, the meaning of physical motion is ``more reluctant to accept the periphrastic form" and ``[d]ouble forms are more typical for figurative meanings than physical ones".

\largerpage[1.5]

Finally, among the ipfv. simplex stems for which the expression of future with the \textit{po}-prefix is considered marginal (against a dominant \textit{bud}-future) we find quite a few that are also attested in the infinitive and other forms of the regular paradigm (e.g., Cz. \textit{vládnout} ‘rule, have control over’, \textit{růst} ‘grow’,\textit{ kvést} ‘blossom’; Saicová Římalová, p.c.). Similarly, we find \textit{po}-prefixed infinitives, for instance \textit{postáti }‘stand’ and\textit{ poležeti} ‘lie’, in Old Czech (\cite{Bondarko1961}: 529, fn. \ref{fn:9}); for these, however, the question arises whether they were not better characterized as pfv. (see Section \ref{sec:wiemer:4.3}). Now that our survey has brought us to this question, we may turn to the diachronic situation.

\subsection{Diachronic background}\label{sec:wiemer:4.2}

There have been different approaches to the ipfv. \textit{po}-future in Czech (and some other Slavic languages) from a diachronic perspective (for partial surveys cf. \cite{Mrhačová1993}: 60--64, \cite{Bláha2008}: 26--34). They have different implications for the question how well established the PFV:IPFV opposition already had been by the 14\textsuperscript{th}--15\textsuperscript{th} centuries, and how much this hinged on the prefixation of verb stems. For instance, according to \citet{LamprechtBauer1986}, in Old Czech certain prefixes with an ingressive meaning (\textit{vz}-, \textit{po}-),\footnote{For examples with \textit{vz}- in Old Czech cf. \citet[192]{Trávníček1923}.} when attached to stems with present tense inflection, gradually turned their ingressive feature into a future meaning (\textit{begin} [to stem\textsubscript{prs}] > \textit{will} [stem\textsubscript{prs}] = \textsc{fut} [stem]). Under this premise, the future function of these prefixes was not dependent on a pre-existing perfectivity opposition (marked by an opposition of simplex vs prefixed stems).

Most scholars, however, base their assumptions (mostly implicitly) on the premise that the perfectivity opposition had already been stable to a sufficient extent in Old Czech times (i.e. by the 14\textsuperscript{th} century). Thus, \citet{Bondarko1961} shows that forms like \textit{půjdu }‘I (will) go’, \textit{ponesu }‘I (will) carry’ in Old Czech behaved like pfv. stems, and, till the end of the 14\textsuperscript{th} century, these stems were also encountered in various types of habitual or timeless (“gnomic”) present:

\ea \label{ex:wiemer:29}
    \textit{Člověk vždy na vše séhne}\textsuperscript{\textsc{pfv.prs}}, \textit{dobrého spieše \textbf{poběhne}}\textsuperscript{\textsc{pfv.prs}}, \textit{ve zlem jsá, pak sě nehne}.\\
    ‘Man always reaches for everything, he rather \textbf{runs} from [= avoids] the good, being unright he nonetheless doesn’t bow.’\\
    \hfill (Alx. V 160\textsuperscript{db}, 43)
\z

\ea \label{ex:wiemer:30}
    \textit{Sto jich jednú se zdi zabijí}\textsuperscript{\textsc{pfv.prs}}, \textit{a však sě vždy ke zdi hrabí}; \textit{pobitých dvě stě \textbf{povlekú}}\textsuperscript{\textsc{pfv.prs}}, \textit{avšak zavše pod zed \textbf{potekú}}\textsuperscript{\textsc{pfv.prs}}.\\
    ‘He knocks off 101 of them from the wall, and yet he always rummages around the wall; \textbf{I carry away} 200 of the knocked ones, but always \textbf{I fall} (back) under the wall.’\\
    \hfill (Alx. V 167\textsuperscript{bß}, 68; \citealt{Bondarko1961}: 532) 
\z

% \largerpage[1.5]

Moreover, \citet[533]{Bondarko1961} argues that all alleged examples of “durative” (i.e. progressive) use also allow for an ingressive reading; he adds the caveat that often the available contexts do not allow for disambiguation among these readings \citep[534]{Bondarko1961}.\footnote{Cf. also \citet[274--275]{Němec2009b}, who however argues for a progressive (“durative”) reading.} Bondarko, however, does not consider the possibility of dispositional readings. Thus, compare \xref{ex:wiemer:31} from the late 14\textsuperscript{th} century, in which the adverbial \textit{bez ustánie} ‘without interruption’ seems to force a progressive reading, with \xref{ex:wiemer:32}, which allows for either progressive or ingressive reading – in either case regardless of deictic time reference. In other contexts, \textit{po}-prefixed verbs of directed motion take part in describing future-oriented single events; see \xref{ex:wiemer:33}:

\ea \label{ex:wiemer:31} 
    \textit{svatý kříž, jenž jest dřevo života těm všěm, ješto jej přijímají, že \textbf{pobyéhnu}}\textsuperscript{\textsc{pfv.prs}} \textit{vesele a \textbf{poneʃʃu}}\textsuperscript{\textsc{pfv.prs}} \textit{bez ustánie kříž svój.}\\
    ‘The holy cross, which is the tree of life for everybody who accepts it, \textbf{I (will) run} cheerfully and \textbf{carry} my cross without interruption.’\\
    \hfill (Matoušova evangelie, 287\textsuperscript{b}.\textsuperscript{14}; \citealt{Bondarko1961}: 533)
\z

\ea \label{ex:wiemer:32} 
    \textit{do města svého... nemohu přijíti, leč \textbf{se} tak \textbf{povezu}}\textsuperscript{\textsc{pfv.prs}} \textit{se pánem Ježíšem.}\\
    ‘I cannot come to my place, but \textbf{I (will) take a ride} with my Lord Jesus.’\\
    \hfill (ChelčPost 245b; \citealt{Němec2009b}: 274)
\z

\ea \label{ex:wiemer:33} 
    \textit{\textbf{Pójdem}}\textsuperscript{\textsc{pfv.prs}} \textit{pod tuto horu, dětem a skotu odpočinem}\textsuperscript{\textsc{pfv.prs}} \textit{a snad sě tu s túhú minem}\textsuperscript{\textsc{pfv.prs}}.\\
    ‘\textbf{I will go} under this hill, I will give a rest to the children and the cattle, and perhaps here I will pass away with longing.’\\
    \hfill (Dal. L 3a, 20; \citealt{Bondarko1961}: 533)
\z

In turn, \citet[440, fn. \ref{fn:14}]{Machek1962} considers \xref{ex:wiemer:31} to exemplify a dispositional reading. In fact, divergent interpretations of these (and many other) examples often cannot be disambiguated by the available context. One may even suspect that, at least in some cases, such “rivalling” readings are more or less equally likely since the \textit{po}-prefixed present tense forms, allowing for any of these readings, oscillated between them.

\subsubsection{Comparison with Latin} \label{sec:wiemer:4.2.1}

A comparison of \textit{po}-prefixed present tense forms as translational equivalents of the Latin future does not clarify the actional character, either: in many cases the ingressive semantics of \textit{po}- + present tense translates the Latin future of \textit{coepio} ‘I begin’, e.g. \textit{puojde} as equivalent of \textit{coeperit}\textsuperscript{\textsc{fut}}\textit{ proficisci} ‘he will (start to) go (against)’. However, we also find \textit{po}- + present tense as a translation of Latin future or present unspecified for actionality. For instance, \textbf{\textit{podme}}\textit{ duch jeho a \textbf{potekú} vody} ‘his spirit \textbf{will blow} and the waters \textbf{will flow}’ (Psalm 147) translating Lat. \textit{flabit}\textsuperscript{ \textsc{fut}}\textit{ …fluent}\textsuperscript{ \textsc{prs}}, or \textbf{\textit{pomyšl’u}}\textit{ za hřiech svój} ‘\textbf{I will consider }as my sin’ for Lat. \textit{cogitabo}\textsuperscript{\textsc{fut}} ‘I will think, consider’ may be interpreted either as processes, or states, or as their incipient phases. On this basis, \citet[78--79]{Šlosar1981} concludes that this behaviour ambiguous to (or oscillating as for?) actionality precludes a clear decision as for the aspect membership of the Old Czech stems (\textit{po-dúti} ‘blow’, \textit{po-téci} ‘flow’, \textit{po-mysliti} ‘think’); see also \citet[102]{Šlosar1981} with regard to the prefix \textit{vz(e)}-.


\subsubsection{Comparison with Old Church Slavonic} \label{sec:wiemer:4.2.2}
\textit{Po}- with motion verbs in the function of marking future (but not pfv. aspect) has already been pointed out for OCS, e.g. by \citet[274]{Němec2009b}, who provides the following example:\footnote{More examples are discussed in \citet[283, 299]{Dostál1954}.}


\ea \label{ex:wiemer:34} Old Church Slavonic\\
    \textit{i milostь tvoě \textbf{poženetъ} mję vьsję dьni života moego.}\\
    ‘And your love \textbf{drives / will drive} me every day of my life.’\\
    \hfill (Ps22:6)
\z

However, again, many (if not most of the) examples adduced in the literature need not unequivocally be interpreted as marking a process in the future, they may also be read as ingressives.\footnote{On the ingressive function of prefixes in OCS and Old East Slavic cf. \citet{eckhoff_2017}.} \citet[77]{Šlosar1981} claims that \textit{po}- as marker of ipfv. future in Old Czech increased in relation to OCS; however, he also adduces examples from OCS sources, e.g. with activity verbs, which correspond to \textit{bud}- forms in later sources of Old Czech. Compare, for instance, \textit{sěmę moě \textbf{porabotajetъ} jemu} ‘my kin/tribe \textbf{will work} for him’, \textit{poleštǫ} ‘I will fly’ in OCS, for which we find \textit{bude slúžiti} ‘will serve’ and \textit{létati budu }‘I will fly’, respectively, in Old Czech sources (\cite{Šlosar1981}: 78). We may, thus, suspect that, apart from a proper account of actionality, we have to distinguish between the (token) frequency of \textit{po}-prefixed stems of (presumably) ipfv. aspect on text level and the productivity in relation to verb stems. Especially the latter appears to have been quite unsteady.



\subsection{Parallel changes concerning past tense forms}\label{sec:wiemer:4.3}

As pointed out by \citet[190--191]{Trávníček1923}, \citet[529, fn. \ref{fn:9}]{Bondarko1961}, and \citet{Šlosar1981} for Old Czech, ipfv. stems which could mark future by prefixing their present tense forms with \textit{po}- (or some other prefix) were also attested with the same prefix in infinitives and other paradigmatic forms. This obviously changed after the 14\textsuperscript{th} century, when these forms, with the exclusion of the present tense and the imperative, gradually disappeared; by the 16\textsuperscript{th} century they became exceptional (\cite{Bondarko1961}: 532--533). Concomitantly, \textit{po}- eventually lost not only its ingressive (i.e. temporal) function, but also its original spatial meanings (‘on a surface’ or ‘from’ = ablative meaning; see \ref{ex:wiemer:29}),\footnote{On the original meanings of \textit{po}- cf. \citeauthor{Němec1958} (\citeyear{Němec1958}: 62--63; \citeyear{Němec2009a}), also \cite{Šlosar1981}: 75--77.} whereas the imperative forms survived because they retained their semantic specialization (see fn. \ref{fn:35}). All this together caused the paradigmatic isolation of the \textit{po}-prefixed present tense forms, which could then be more easily associated with the paradigm of the respective simplex stems belonging to ipfv. aspect and marking activities or states.

The loss of the ingressive function was an important factor contributing to the paradigmatic isolation of the \textit{po}-prefixed present tense forms, and this loss (or lack of further development) should also be considered in connection with the past tense. Past tense forms of the same stems, but without a prefix, are well attested in contexts in which they gained ingressive meaning; for instance 

\ea \label{ex:wiemer:35} Old Czech\\
    \textit{a vstav \textbf{šel} jest po \textbf{něm}.}\\
    ‘He then stood up and \textbf{went} after him / \textbf{followed} him.’\\
    \hfill (BiblLit Mc2:14; cited from \cite{Němec2009b}: 277)
\z

That is, the ingressive phase of some activity was usually only induced from the context, but not marked by a prefix that also would have enhanced the chance of this stem being classified as a pfv. one.

\citet[542--543]{Bondarko1961} shows that lacking prefixation in the case of ingressivity was typical of Old Czech (for both past and non-past forms), that is, simplex stems also clearly prevailed over their prefixed derivatives (which could then easily be interpreted as pfv. stems) in this function. See \xref{ex:wiemer:36} for a sequence of events (\textit{vzvědě}.\textsc{aor.3sg} ‘learnt’ > \textit{jěde}.\textsc{prs.3sg} ‘goes, rides’) and \xref{ex:wiemer:37} with a contrast between \textit{pojedu} and \textit{jěde}: both forms refer to the initiation of movement, but the prefixed form in the quote (\textit{pojedu}.\textsc{1sg}) can easily be understood as a commissive speech act. It can thereby evoke an immediate future, whereas the latter form (\textit{jěde}.\textsc{3sg}) pairs up with an (already obsolete) aorist-form (\textit{poče}.\textsc{aor.3sg}) to mark a narrative sequence.

\ea \label{ex:wiemer:36} Old Czech\\
    \textit{Ciesař Bedřich poslední sňem zapovědě. Jakž to král český \textbf{vzvědě}, k dvoru \textbf{jěde}.}\\
    ‘Emperor Friedrich announced the latter in the assembly. As the Czech king \textbf{learnt} about this, he went [lit. \textbf{goes}] to the court.’\\
    \hfill (Dal 75\textsuperscript{a}, 139)
\z

\ea \label{ex:wiemer:37} Old Czech\\
    \textit{„\textbf{Pojedu} já proti otci svému, snad to přivedu k konci dobrému.“ Tak i \textbf{jěde} proti otci svému i \textbf{poče} dobývati milosti svému milému.}\\
   ‘“\textbf{I (will) go} against my father, maybe I will lead this to a good end.” So he went [lit. \textbf{goes}] against his father and \textbf{began} to earn love for his dear.’\\
    \hfill (Dal 38\textsuperscript{b}, 79; cited from \cite{Bondarko1961}: 537)
\z

Notably, simplex stems denoting directed motion – as well as other simplex stems denoting activities or states – usually did not mark ingressivity by prefixes in earlier stages of other Slavic languages, either (cf. \cite{Bondarko1961}: 537--538, \citealt{Dickey2011}: 185--186).\footnote{Only in the 17\textsuperscript{th} century did the ingressive and other phases of atelic activities start being marked more and more consistently by (pfv.) prefixed derivatives in the eastern half of Slavic and in Polish (\cite{Dickey2011}).}

In general, the integration of atelic stems (and of stems denoting punctual events) into the aspect system (via consistent derivation) was belated by a couple of centuries even in those Slavic languages that can be considered “most progressive” in this respect, i.e. in the eastern half of Slavic (cf. \cite{Bermel1997} for East Slavic). In fact, after the 16\textsuperscript{th} century the eastern part of Slavic (incl. Polish and Bulgarian) considerably extended the use of \textit{po}- as a productive device for deriving delimitatives (= pfv. activities) and as an ingressive marker of directed motion.\footnote{Cf. \citeauthor{Dickey2011} (\citeyear{Dickey2011}; \citeyear{dickey_2018}: 375--376). In her study based on the TOROT treebank, \citet[]{eckhoff_2020} found that prefixation of atelic ipfv. stems with \textit{po}- had already been quite productive in OCS and Old East Slavic. In other words, prefixation as a means of perfectivizing simplex stems had already advanced to an extent that this derivational pattern could occur without a concomitant telicizing effect (see Section \ref{sec:wiemer:3.1.1}). However, it obviously needed some more centuries for this pattern to become productive for the derivation of pfv. stems with delimitative meanings (compare Russ. \textit{po-sidet’} ‘sit a while’).} After all, Czech and other languages on the western edges of the Slavic speaking territory are remarkably different from Slavic languages in the eastern half (East Slavic, Bulgarian, but also Polish). In the eastern half, ingressivity in (narrative) past contexts is consistently marked with prefixed pfv. equivalents of the respective ipfv. stems, whereas in the western group this explicit marking is largely absent (cf. \citealt{Ivančev1961}, \citealt{Berger2016}). Even more remarkably, the areal range of this latter group is more or less coextensive with the use of prefixed ipfv. verbs to mark future, as has been noted for Slovak, Sorbian and (although to a lesser degree) for Slovene and SerBoCroatian (cf. \citealt{Dickey2011}: 178--180).


\subsection{Summary of the presumable scenario}\label{sec:wiemer:4.4}


Table \ref{tab:wiemer:2} schematically summarizes the presumable diachronic changes that affected \textit{po}-prefixed present tense forms of certain semantic verb groups (with verbs denoting directed motion as the core). It basically reproduces the view expounded in \citet[277]{Němec2009b}. His “story” can be retold in the following way: in Old Czech the aspect opposition (PFV:IPFV) had already been well established, and prefixes served to derive pfv. stems from ipfv. simplex stems. However, verbs of motion created a problem since their non-past (“present”) tense forms occurred in contexts in which the prefix could be interpreted as marking either ingressivity or duration (i.e. progressive or, rarely, stative meaning); in certain contexts either interpretation was possible. Ingressive readings favoured the categorization of the prefixed form as representing pfv. aspect, whereas progressive readings favoured its classification as ipfv. aspect. In this regard, the prefixed non-past forms of these stems were underspecified. Simultaneously, up to a certain time, the same prefixed stems also occurred in the past tense, but, contrary to the non-past forms, the prefix here clearly marked ingressivity, so that the stem could easily be assigned pfv. aspect. With the simplex stem, in turn, actionality was underspecified, as it was with the non-past forms of the respective prefixed stem. This led to paradigmatic asymmetry. Subsequently, this asymmetry was “reshuffled” in favour of treating the prefixed non-past forms as ipfv. (with future reference), probably because unambiguously ingressive contexts proved to be in the minority. Concomitantly, the prefixed past tense forms of these verbs vanished (unless in other than motion meanings), so that the prefixed non-past forms could be integrated into the tense paradigm of the unprefixed forms (for past, \textit{šel jsem }‘I went’, and present, \textit{jdu }‘I am going’) as their future. In this way, the biaspectual behaviour of \textit{pójdu} > \textit{půjdu} was abandoned and \textit{jít} ‘go’ became an \textit{imperfectivum tantum}. Table \ref{tab:wiemer:2} sketches these changes with the 1sg-forms of \textit{jít(i)}‘go’, which was the most frequent motion verb in Old Czech.

\begin{table}
\caption{Reinterpretation of paradigmatically isolated \textit{po}-prefixed present tense}
\label{tab:wiemer:2}
\begin{tabularx}{\textwidth}{p{3.1cm}Xp{1.9cm}p{3cm}}
\lsptoprule
\textit{jít}(\textit{i}) ‘go’ & {Past} & {Present} & {Future} \\
\midrule
Old Czech & PFV: \textit{pošel jsem} & \raisebox{3\height}{\rule{0.8cm}{0.25mm}} & PFV/IPFV: \textit{pójdu}\footnote{\label{fn:a} Periphrastic *\textit{budu jíti} is unattested and probably never existed.}\\
(14\textsuperscript{th}--15\textsuperscript{th} century) & IPFV: \textit{šel jsem} & IPFV: \textit{jdu} & \\
\midrule
% \multirow{2}{*}
{Intermediate stage} & PFV: \textit{pošel jsem} & \raisebox{3\height}{\rule{0.8cm}{0.25mm}} & \raisebox{3\height}{\rule{0.8cm}{0.25mm}}\\ 
& IPFV: \textit{šel jsem} & IPFV: \textit{jdu} & PFV/IPFV: \textit{pójdu}\\
\midrule
Modern Czech & IPFV: \textit{šel jsem} & IPFV: \textit{jdu} & IPFV: \textit{půjdu} \\
\lspbottomrule
\end{tabularx}
\end{table}


In addition to \citeauthor{Němec1958}’s considerations, one has to account for the rise of \textit{bud}- as a future auxiliary with ipfv. stems, and as such \textit{bud}- (+ infinitive) started interfering with \textit{po}-prefixed ipfv. stems, although quite unevenly, depending on the verb lexeme (see Section \ref{sec:wiemer:4.1}). To a large extent, \citeauthor{Němec1958}’s argument is based on \textit{jít(i)} ‘go’ and the absence of periphrastic *\textit{bud- jít(i)} in the history of Czech (see fn. \ref{fn:a}). This fact precluded any potential morphological opposition (\textit{pójdu} vs *\textit{budu jíti}) just for this particular verb (with direct motion meaning); however, for many other items such an opposition was possible, since “doublets” of \textit{po}-prefixation and \textit{bud}- + infinitive already existed in Old Czech. Although \citet[79--80, 102--103]{Šlosar1981}, on the basis of the old text corpus and older grammars, argues that the number of ipfv. verbs expressing their future with \textit{po}- (or another prefix) had already diminished in older periods of Czech (see Section \ref{sec:wiemer:4.2.2}), other findings summarized in Section \ref{sec:wiemer:4.1} suggest that the duality of future marked by present tense with \textit{po}- vs \textit{bud}- (+ infinitive) is evident for a quite large and open class of verbs and that it has probably been pervading the entire history of Czech.

In sum, as concerns the approaches referred to in Section \ref{sec:wiemer:4.1}, most researchers assume (often implicitly) that the perfectivity opposition must already have been stable to a sufficient extent by the 14\textsuperscript{th} century.\footnote{The argument in \citet[]{Němec2009b} does not contradict \citegen{Bondarko1961} argument. After all, the considerations of both scholars are based on the idea that the \textit{po}-prefixed forms lost their ingressive function and that the isolated \textit{po}-prefixed present tense forms were integrated for symmetry reasons into the paradigm of the ipfv. stem.} Regardless of this assumption (and how it was argued for), all scholars appear to agree that a former ingressive meaning of the relevant prefixes turned into a function of tense, which is by itself insensitive to aspect. Whether the compatibility with progressive contexts (e.g., with durative time adverbials) is an ancient feature (connected to the individual “lexical prehistory” of the prefixes) or only became possible after these prefixes were reinterpreted as marking future tense of ipfv. simplex stems, is of secondary concern.



\section{Chronological relations and areal distribution}\label{sec:wiemer:5}

The differentiation of emergent future auxiliaries in various regions of Slavic is already discernible at the time of the first written documents after early OCS, i.e. from the 11\textsuperscript{th} century onward. In the eastern part of South Slavic (= Balkan Slavic), \textsc{want} superseded its strongest “rival” \textsc{have}, although \textsc{have}-based patterns have remained important until today; \textsc{become} did not play any role, quite to the contrary of the western part of Slavic, both in the North and the South. In turn, the eastern part of North Slavic (= East Slavic) for a couple of centuries showed preference for \textsc{begin}-based periphrases (with \textsc{want}-based patterns probably imported from OCS), before \textsc{become} spread there as well and became the dominant or exclusive auxiliary of futures with ipfv. stems.

The employment of prefixes to mark the future of ipfv. simplex stems (and not to derive a pfv. stem) looks quite exotic, also from a Slavic point of view. If the reconstruction sketched in Section \ref{sec:wiemer:4} is correct, this could mean two things. Either, by the 14\textsuperscript{th} century, a stem\hyp{}derivational opposition of pfv. vs ipfv. aspect had already become sufficiently stable in general, but then in the western parts of the Slavic-speaking territory this system turned out less consistent for some groups of verb lexemes that denoted parametric changes or atelic activities. Such an interpretation is called for if we assume a reinterpretation of the relevant prefixes in the paradigm of the relevant ipfv. stems (see Section \ref{sec:wiemer:4.4}). Or, alternatively, we may argue that the use of \textit{po}- and other prefixes as markers of ipfv. future reflects an archaic trait which, as it were, survived the first stages in the development of the stem\hyp{}derivational aspect system and which also manifested itself earlier in OCS\footnote{That \textit{po}- for ipfv. future might be an archaism was assumed by \citet{Němec2009b} and by \citet[48]{Galton1976}. However, Němec also at least implicitly advocated for ``paradigmatic reshuffling'', which is inherent to the first alternative (see the previous footnote).} (see Section \ref{sec:wiemer:4.2.2}). 

At present, we cannot decide between these alternative explanations, but they do not even exclude one another. Verb stems denoting unbounded activities represent actionality groups that are more “vulnerable” within a perfectivity system whose onset and initial stages were certainly triggered by distinctions related to telicity; into such a system lexemes denoting unbounded activities were integrated only at a later stage (see Section \ref{sec:wiemer:3.1}). Their integration could fade away, or get stuck at an early stage, and this is what evidently happened in Czech, as part of a “cluster” of languages of the “Slavic west” – in contrast to the eastern part, in which prefixes are much more freely employed for the derivation of pfv. delimitatives and ingressives of activity verbs. Scholars like \citeauthor{Dickey2011} have made a plausible case that contact with German was the most likely factor that weakened and hampered the use of prefixes for the derivation of pfv. stems with lexemes denoting unbounded activities (see Section 4.3). \citet[]{Dickey2011} calls this ‘replica preservation’. What was preserved (or, maybe, rather restored) was the lack of an explicit differentiation between temporally limited and unlimited activities by otherwise well-developed derivational means (i.e. prefixes). However, in light of data discussed in Section \ref{sec:wiemer:4.2}--\ref{sec:wiemer:4.3}, this “preservation” was obviously preceded by a period (14\textsuperscript{th}--15\textsuperscript{th} centuries) of rather free alternation between simplex (ipfv.) and prefixed (pfv.) stems for ingressivity (both in past and non-past tense) as well as for future-related and habitual or gnomic situations (with non-deontic modal implicatures).

\largerpage[1.5]

When we now resume the “story” about the \textsc{become}-future, more precisely: about the pattern \textsc{become} + infinitive, in North Slavic, we notice that the areas from which this pattern most probably spread overlap strikingly with the areal cluster in which prefixation had been showing some perceivable tendency to mark the future of ipfv. stems. More so, even the period of these processes largely coincides: we are dealing here with the 14\textsuperscript{th}--16\textsuperscript{th} centuries. The innovative pattern \textit{bud}- + infinitive became prominent first in Bohemian (Czech), then in Polish, from where it spread into western East Slavic (Ruthenian) and finally into Russian (15\textsuperscript{th}--18\textsuperscript{th} century; see also the map in \cite{Whaley2000}: 74). Properly speaking, this spread from Bohemia toward the northeast meant that a minor pattern\footnote{Such patterns were also assumed by \citet[30]{andersen_2006b} for prehistoric dialect zones all over from today's Serbia to Danzig. He however also admits that this scattered spread remains mysterious.} was supported and amplified by inner-Slavic contact from the (south)west of North Slavic, and that it was mainly characteristic of written (pre-standard) varieties. This interpretation is also consistent with the observation that North Slavic belongs to a larger area stretching from German to northeastern Europe in which futures are either lacking or are only weakly developed on the basis of \textsc{become} (\citealt{wiemer_2020}: 276, following \citealt{Dahl2000}). This makes \textsc{become} areally salient as a source of future auxiliaries. Contact with Middle and Early New High German can be considered a plausible catalyst of the gradual entrenchment of inchoative \textit{bǫd}-/\textit{bud}- in Bohemia and, thus, in all North Slavic, although this issue continues to create some problems, first of all because of the provenance of the infinitive (see Section \ref{sec:wiemer:3.2.1.3}). 

Moreover, the innovative pattern \textit{bud}- + infinitive is also attested in South Slavic, but only in its “middle section” (i.e. Štokavian-Čakavian). Within that region, it spread from west to east, more or less at the same time when in the eastern part the \textsc{want}-future developed and spread from east to west, concomitantly with the stepwise loss of the infinitive. The pattern \textit{bud}- + infinitive has never been attested in Slovene, i.e. the northwest periphery, for which influence of German (apart from Italian) varieties was certainly more likely than on the Čakavian Adriatic coast. We cannot know whether this lack of attestation in Slovene simply arises from lack of data about its earliest manifestations in texts.

Anyway, the distribution of \textit{bud}- + infinitive in South Slavic looks strange: its area of use is separated from German-speaking territory by Slovene, which, in turn, preserves the earlier pattern \textit{bô}- (< *\textit{bǫd}-) + \textit{l}-form (simultaneously losing its future perfect function) and is isolated in this respect. One wonders whether there might have been a stable connection to West Slavic, namely Bohemia and Moravia despite the separation by the Magyars (since approx. 900 AD). This would not explain the preservation of the isolated \textit{bud}- + \textit{l}-form pattern in Slovene, but it would fit the western cluster of Slavic languages (Czech, Slovak, Sorbian, Slovene), which have shown some tendency to use prefixes as markers of the future with ipfv. stems. This cluster corresponds to an earlier continuum between what subsequently became Bohemian, Moravian and Kajkavian dialect zones; this cluster must have existed prior to the arrival of the Magyars in the Pannonian plain, and to the exclusion of the Lekhitic territory farther to the northeast (from where Polish and Kashubian developed). This cluster was also pointed out by \citet[25]{andersen_2006b}, but without an account of prefixes as markers of future with ipfv. simplex stems. One thus wonders whether the parallelism of these two processes was simply a coincidence. What can be said with more confidence is that, to a large extent, different dominant future markers (\textsc{become vs want}) developed due to divergent contact-induced processes in different parts of the Slavic-speaking area and of inner-Slavic continua after that time.

How does all this relate to the hypoanalysis of PFV.PRS to pfv. future in North Slavic, and to its non-occurrence in South Slavic? First of all, since this kind of hypoanalysis was based on the stem\hyp{}derivational aspect opposition, this opposition must have been strong enough by the time when among looser combinations with lexical source expressions (\textsc{want, become, have, begin, take}) areally biased preferences started acquiring shape. 

Since Common Slavic had no dedicated contrasts in the non-past domain, future reference had been a natural part of the meaning range of non-past (“present”) tense, regardless of aspect, prior to the emergence of futures based on auxiliaries. However, PFV.PRS was driven into the irrealis domain (“inactual present”) by the aspect (IPFV:PFV) opposition alone. Subsequently, future readings must have increasingly outnumbered irrealis readings in North Slavic, but irrealis readings did not vanish; instead, this usage domain of PFV.PRS is still shared with South Slavic. Here, in turn, the rise of dedicated future auxiliaries also introduced a morphological contrast with pfv. stems, which is alien to North Slavic; however, remnants of genuine future uses of PFV.PRS can be found especially in the northwestern part of South Slavic to this day (cf. \citealt{Grković-Mejdžor2019}: 11). 

Consequently, irrealis uses have remained the large domain in which PFV.PRS continues to appear both in North and South Slavic. In this respect, nothing has changed since CS times: the non-future uses, considered “residual” for North Slavic, are just as original as are uses of PFV.PRS to refer to future events. What obviously has changed is the salience of future reference for PFV.PRS. Therefore, an explanation is required for (at least) two issues: first, how can this shift in salience be captured, also within a more overarching framework of language change? Second, should the divergent development of PFV.PRS in North and in South Slavic be explained by differences in the development of the stem\hyp{}derivational aspect system?

While the first question requires another methodological take on the topic, there is no reason to give an affirmative answer to the second question. For all we know about inner-Slavic differences in the consistency with which stem\hyp{}derivational aspect has been developing, the eastern half appears to be more consistent in marking aspect with stems of all actionality classes, i.e. lexical expansion is more advanced, and the respective process seems to have begun more or less in the 16\textsuperscript{th} century (see Section \ref{sec:wiemer:4.3}). After all, we identified the use of prefixes marking the future with ipfv. stems that denote activities and parametric changes as conditioned by a system that existed prior to this process. As for the centuries preceding this process, there is no reason to assume any remarkable differences in the development of the aspect system, either, despite the fact that a restriction of \textsc{become} + infinitive to ipfv. stems applying from the first documents has been reported for North Slavic but not for South Slavic languages (see Section \ref{sec:wiemer:3.2.1}). Moreover, we are left with the finding that the “future-by-prefixes” pattern, jointly with the initial spread zones of \textsc{become} + infinitive, cluster in the west of Slavic and cut the North-South divide. That is, they are orthogonal with regard to the North-South split. The lack of aspect restrictions for the \textsc{want}-future can be explained rather straightforwardly from the source construction, but the restriction to ipfv. aspect for the \textsc{become}-future in the North cannot. Nor can this restriction be explained by the most plausible “external catalyst” of \textsc{become} + infinitive in West Slavic, namely German.

\section{Conclusions}\label{sec:wiemer:6}

The study of the development of futures in Slavic languages shows that a clear-cut distinction between grammaticalization and hypoanalysis as mechanisms of change not only makes sense, but that such a distinction is even necessary in order to see their entanglement and relative chronology as well as the divergent consequences this has for the formation of areal clusters and clines within the whole language group. Croft's notion of hypoanalysis has here been adopted with two slight, but important amendments for the case under study: first, we need to apply it to paradigmatic oppositions, regardless of their morphosyntactic realization; second, the default future function of PFV.PRS in North Slavic results not so much from a shift towards a new function that previously was only implied (i.e. future), but this function was always possible, before it turned into the most prominent one. It remains to model this mechanism within a more comprehensive framework of linguistic changes. 

Concomitantly, we should correctly describe the stem\hyp{}derivational nature of the Slavic perfectivity opposition in order to understand the fundamental role which this technique plays in the grammatical outfit of Slavic languages in general, and for the verbal system in the non-past domain in particular. With this in mind, PFV.PRS turns out to be an extremely persistent feature of \textit{all} Slavic languages, inasmuch as PFV.PRS in all of them continues to serve the same functions in the irrealis domain as it most probably has done from Common Slavic times, when the perfectivity opposition had established itself firmly enough. Thus, in this sense, its future function is just an areally restricted by-product that has not ousted other functions.

Finally, the North-South split of future marking leaves open issues. First of all, there is still no satisfying explanation for why the pattern \textsc{become} + infinitive is restricted to ipfv. stems in North Slavic, and why this obviously has been so through the entire attested period, i.e. since the 13\textsuperscript{th} century. Simultaneously, we cannot say with certainty where, when and how the infinitive entered this pattern, in particular whether contact with German varieties provided the decisive external trigger. However, we observe that the initial spread zones of \textsc{become} + infinitive perspicuously overlap with a “western cluster” of those Slavic languages in which prefixes have been used to mark the future of ipfv. simplex stems. I hope to have supplied sufficient arguments that this overlap is not just incidental but testifies to a common period (13\textsuperscript{th}--16\textsuperscript{th} centuries) during which these seemingly unrelated patterns of future marking of ipfv. stems developed. In the western cluster, the \textsc{become}-based future suppressed the prefixation pattern, but at least in Czech is has never ceased to be productive for the relevant actionality classes.


% \section*{Abbreviations}
% \begin{tabularx}{.5\textwidth}{@{}lQ@{}}
% ... & \\
% ... & \\
% \end{tabularx}%
% \begin{tabularx}{.5\textwidth}{@{}lQ@{}}
% ... & \\
% ... & \\
% \end{tabularx}





\section*{Acknowledgements}
I want to cordially thank, first of all, Jana Kocková, Lucie Saicová Římalová and Jasmina Grković-Mejdžor, but also Tilman Berger, Marketa Janebová, Michaela Martinková, Karolína Skwarska and Mladen Uhlik for their expertise and support. I also appreciate the useful remarks by an anonymous reviewer, first of all concerning Central South Slavic and some threads of argumentation. Of course, the usual disclaimers apply.

%\section*{Contributions}
%John Doe contributed to conceptualization, methodology, and validation. 
%Jane Doe contributed to writing of the original draft, review, and editing.

\sloppy
\printbibliography[heading=subbibliography,notkeyword=this]
\end{document}




